\documentclass[output=paper]{langscibook}
\ChapterDOI{10.5281/zenodo.22880230}
\title{The basic linguistic features of the Teop language}
\author{Ulrike Mosel}
\abstract{\noabstract}
\begin{document}
\maketitle\label{chap:3}
\section{Introduction}
Teop is a configurational language. The building blocks of clauses are eight types of phrases which are formed by content words and specific functional words in a strict word order\footnote{The terms of Teop formal categories of words, phrases and clauses are distinguished from terms for typological comparative concepts by capital letters, see \textcite{Haspelmath:2016mf}, e.g. topical subject Noun Phrase.}, see (\ref{ex:sketch-1}). 
\ea\label{ex:sketch-1}
    A ba suburaara na kuukuhi overe ni roho bona kehaa.\\
    \gll a ba subu-raara na kuu{$\sim$}kuhi overe ni roho bona kehaa\\
    [\textsc{art}2.\textsc{sg}	\textsc{kin}.\textsc{pl}	grandparent-1\textsc{pl}.\textsc{in}.\textsc{poss}.\textsc{pron}]\textsubscript{NP} [\textsc{tm}2 \textsc{red}{$\sim$}grate coconut \textsc{appl} before]\textsubscript{VC} [\textsc{obj}.\textsc{art}2.\textsc{sg} shell]\textsubscript{NP}\\
    \glt `Our grandparents grated coconuts with a shell.' (\ref{chap:5}:\ref{ex:text01-75})
\z
The eight types of phrases are distinguished in the following way:
\begin{enumerate}
    \item Noun Phrase (NP): The NP is introduced by an obligatory Article signifying the lexical class of the NP head, which is distinct from its word class, and the specificity of its reference, its grammatical number, and its syntactic function. The head of the NP may be a Noun (N), a Verb (V), or an Adjective (Adj).
    \item Locative Phrase (LocP): The LocP has the same structure as the NP, but lacks an Article. It is headed by a Noun that refers to a specific place.
    \item Verb Complex (VC): The first position of the VC is optionally held by one of the eight particles that signify temporal relations and modality (Gloss: \textsc{tm}1--\textsc{tm}8) and its last position by the Imperfective or the Perfective Aspect Marker (Gloss: \textsc{ipfv}, \textsc{pfv}). Affirmative VCs are headed by a Verb (V), a Noun (N), an Adjective (Adj), or a Cardinal Numeral (Num). In addition, the core of a VC that is negated by the Double Negator may consist of a Noun Phrase (NP), a Prepositional Phrase (PrepP) or a Linker Phrase (LnkP), See \sectref{sec:subphrase} and \sectref{sec:verbcomplintro}.
    \item Adjectival Phrase (AdjP): The AdjP is marked by an Article that formally agrees with the Article of the NP it relates to as a predicate or an attribute. It is headed by an Adjective (Adj) or a Verb (V).
    \item Numeral Phrase (NumP): The NumP is headed by a Cardinal Numeral and marked by an Article that formally agrees with the Article of the NP it relates to in the function of an attribute or anaphoric argument.
    \item Substitute Phrase (SubstP): The SubstP is marked by an Article and headed by \textit{vuaba} `one' or \textit{meha}/\textit{meho} `other' and functions as an anaphoric substitute for a NP. The selection of the Article and the form \textit{meha} or \textit{meho} is determined by its antecedent. 
    \item Prepositional Phrase (PrepP): The PrepPs are introduced by a preposition which is followed by a Pronoun or NP. The most frequent preposition is the Multi-Purpose Preposition \textit{te} `\textsc{prep}'.
    \item Linker Phrase (LnkP): The LnkP is introduced by the Linker \textit{vaa} `\textsc{lnk}' which is followed by an Adverb (Adv), a LocP, a PrepP introduced by \textit{te} `\textsc{prep}', or a complement clause introduced by \textit{tea} `\textsc{compl}1'.
\end{enumerate}
The following grammatical description of Teop starts with syntactically simple Non-Verbal Clauses (\sectref{sec:nonverbal}) and Verbal Clauses (\sectref{sec:verbal}), which is followed by an overview of the syntactic functions of phrases in \sectref{sec:synfunc}, Then the classification of the content words and Numerals is presented in \sectref{sec:numerals} on the basis of their distribution within the phrases. The morphology of content words is described in \sectref{sec:morphology}. The following sections deal with the complex structures of VCs (\sectref{sec:vc}) and NPs (\sectref{sec:np}). The forms and syntactic functions of Pronouns are listed in tables and explained in \sectref{sec:pronouns}. The final two sections are about the Alienable and Inalienable Possessive Constructions (\sectref{sec:posscon}) and dependent clauses (\sectref{sec:depclaus}).

The Teop phonology and orthography have been described by \textcite{Radtke:2005kw}.

Information on inflectional and derivational morphemes, grammatical particles and constructions on phrase level is also given in the Multifunctional Teop-English Dictionary (\textsc{mted}), \textcite{Mosel:2019bt}, which facilitates filtered searches for lexemes, parts of speech, meaning descriptions, bound morphemes, glosses, multi-word constructions, and semantic domains. An introduction to complex filtered search strategies is given in \sectref{sec:mted}.
\section{Non-Verbal Clauses}\label{sec:nonverbal}
There are four types of Non-Verbal Clauses:
\begin{enumerate}
    \item Existential clauses consist of a single NP and state the existence of the NP referent (\ref{ex:sketch-2a}).
    \item In topic-comment clauses, a topical subject NP (\ref{ex:sketch-3}), Pronoun, anaphoric NumP or SubstP (\ref{ex:sketch-2b}, \ref{ex:sketch-6}) is followed by a predicative NP (\ref{ex:sketch-3}), AdjP (\ref{ex:sketch-4}), LocP (\ref{ex:sketch-2b}) or PrepP (\ref{ex:sketch-6}).
    \item In thetic\footnote{\textcite[137--146]{Lambrecht:1994od}, \textcite{Sasse:1987yh}.} non-verbal clauses the non-topical subject form follows the predicate form (\ref{ex:sketch-5}). 
    \item The negative existential clause is formed by the Negator \textit{ahiki} `not exist' and a non-specific NP (\ref{ex:sketch-7}).
\end{enumerate}
The text excerpt below starts with a PrepP referring to the topic which is followed by an existential clause and two topic-comment clauses with predicative LocPs, see (\ref{ex:sketch-2}).

\ea \label{ex:sketch-2}
    \ea \label{ex:sketch-2a}
    Teo potee, o buaku koa o vanai,\\
    \gll te=o	potee	o	buaku	koa	o   vanai\\
    [\textsc{prep}=\textsc{art}3.\textsc{sg}	potee]\textsubscript{PrepP}	[[\textsc{art}3.\textsc{sg}	two	only]\textsubscript{NumP}	\textsc{art}3.\textsc{sg}	outrigger.boom]\textsubscript{NP}\\
    \glt `As for the \textit{potee} canoe, there are only two outrigger booms,'\footnote{NPs with an attributive NumP are always singular.}
    \ex \label{ex:sketch-2b}
    o peho matana, o meho he itao.\\
    \gll 	o	peho	mata-na	o	meho	he	itao\\
    [\textsc{art}3.\textsc{sg}	one3.\textsc{sg}]\textsubscript{NumP}	[front-\textsc{derel}]\textsubscript{LocP}	[\textsc{art}3.\textsc{sg}	other3.\textsc{sg}]\textsubscript{SubstP}	\textsc{advers}	[stern]\textsubscript{LocP}\\
    \glt `one near the bow, and the other near the stern.' (\ref{chap:24}:\ref{ex:text18-44})
    \z    
\z 
In (\ref{ex:sketch-3}) a NP functions as the subject and a second NP as the predicate. The latter is modified by an attributive AdjP, see (\ref{ex:sketch-3}).
\ea\label{ex:sketch-3}
    A goroto a iana a beera sana.\\
    \gll a	goroto	a	iana	a	beera	sana\\
    \textsc{art}2.\textsc{sg}	turtle	\textsc{art}2.\textsc{sg}	swimming.animal	\textsc{art}2.\textsc{sg}	big	very\\
    \glt `The turtle is a very big swimming animal.' (\ref{chap:26}:\ref{ex:text19-2})
\z
AdjPs may also function as a predicate, see (\ref{ex:sketch-4}).
\ea\label{ex:sketch-4}
    O atovo o mataa,\\
    \gll o	atovo	o	mataa\\
    \textsc{art}3.\textsc{sg}	sago.palm	\textsc{art}3.\textsc{sg}	good\\
    \glt `The sago palm is good,' (\ref{chap:6}:\ref{ex:text2-54})
\z
The ambiguity of NP-AdjP constructions is resolved by the context or by \textit{vai} `\textsc{dem}6' which links the attributive AdjP to the NP, see (26) in \sectref{sec:subphrase}.

In thetic clauses, which focus on the whole situation and lack a topical NP, the predicative AdjP precedes the non-topical subject NP, see (\ref{ex:sketch-5}).
\ea\label{ex:sketch-5}
    A pereperee ge kasikasiree a kapa nae.\\
    \gll a	pereperee	ge	kasikasiree	a	kapa	nae\\
    \textsc{art}2.\textsc{sg}	green	or	yellow	\textsc{art}2.\textsc{sg}	cover	3\textsc{sg}.\textsc{poss}.\textsc{pron}\\
    \glt `(When it grows bigger, they call it \textit{magura} `ripening coconut.') Its cover is green or yellow.' (\ref{chap:5}:\ref{ex:text01-38})
\z
A predicative PrepP is found in the description of making a net for catching turtles,\footnote{A predicative PrepP introduced by \textit{mene} occurs in (\ref{chap:16}:\ref{ex:text11-37}).} see (\ref{ex:sketch-6}). 
\ea\label{ex:sketch-6}
    \ea\label{ex:sketch-6a}
    O peho teo meho punete,\\
    \gll o	peho	te=o	meho	punete\\
    [\textsc{art}3.\textsc{sg}	one3.\textsc{sg}]\textsubscript{NumP}	[\textsc{prep}=\textsc{art}3.\textsc{sg}	other3.\textsc{sg}	edge]\textsubscript{PrepP}\\
    \glt `(When this work is finished, we get two helmet shells for the bottom of the net.) One (is) for one edge,'
    \ex\label{ex:sketch-6b}
    o meho he teo meho punete.\\
    \gll o	meho	he	te=o	meho	punete\\
    [\textsc{art}3.\textsc{sg}	other3.\textsc{sg}]\textsubscript{SubstP}	\textsc{advers}	[\textsc{prep}=\textsc{art}3.\textsc{sg}	other3.\textsc{sg}	edge]\textsubscript{PrepP}\\
    \glt `the other (is) one for the other edge.' (\ref{chap:27}:\ref{ex:text20-28})
    \z
\z
Negative existential clauses are formed by the Negator \textit{ahiki} `not exist'\footnote{The Negator \textit{ahiki} `not being there' can be analyzed as a NP formed by the Article \textit{a} `\textsc{art}2.\textsc{sg}' and the Verb \textit{hiki} `not exist, not have', \url{https://dictionaria.clld.org/units/teop-hiki_1}, but it is never written as two words.} and a NP marked by a Non-Specific Article (see \sectref{sec:paraarticles}). The example below comes from a text about a pandanus species, see (\ref{ex:sketch-7}).
\ea\label{ex:sketch-7}
    Ahiki to toutou teve.\\
    \gll ahiki	to	toutou	teve\\
    not.exist	[\textsc{nspec}3.\textsc{sg}	thorns	\textsc{prep}.3.\textsc{sg}.\textsc{pron}]\textsubscript{NP}\\
    \glt `It does not have any thorns.' (lit. `There aren't any thorns of it.') (\ref{chap:13}:\ref{ex:text9-18})
\z
Predicative NPs and PrepPs are negated by the verbal Double Negator Construction, see \sectref{sec:subphrase}. Predicative AdjPs are negated by a VC which is marked by the Double Negator and headed by the AdjP, compare (\ref{ex:sketch-4}) with (\ref{chap:20}:\ref{ex:text15-10}).

\section{Verbal Clauses}\label{sec:verbal}
In topic-comment Verbal Clauses, the VC follows a clause initial topical NP, Pronoun, or anaphoric phrase which in transitive clauses may function as the subject or the object. Topical NPs and non-topical subject NPs are marked by a Basic Article. But non-topical object NPs are only marked by the Basic Article (See \sectref{sec:paraarticles}), if the topical subject is formed by a 1\textsuperscript{st} or 2\textsuperscript{nd} person Pronoun, see (\ref{ex:sketch-8}).
\ea\label{ex:sketch-8}
    Eara repaa karem vakakaniva a peha papana.\\
    \gll e-ara	re= paa	karem	va-kakaniva	a	peha	papana\\
    \textsc{h}-1\textsc{pl}.\textsc{in}.\textsc{pron}	\textsc{consec}= [\textsc{tm}3	scrape	\textsc{caus}-sharp]\textsubscript{VC}	\textsc{art}2.\textsc{sg}	one2.\textsc{sg}	end\\
    \glt `Then we sharpen one end.' (\ref{chap:5}:\ref{ex:text01-55})
\z
But if an NP or a 3\textsuperscript{rd} person Pronoun functions as the topical subject, the object NP is marked by the Object Article, see (\ref{ex:sketch-9}).
\ea\label{ex:sketch-9}
    A abana na mosi maari bono kopihoo ...\\
    \gll a	abana	na	mosi	maa=ri	bono	kopihoo\\
    \textsc{art}2.\textsc{sg}	men	[\textsc{tm}2	cut	hither=3\textsc{pl}.\textsc{ipfv}]\textsubscript{VC}	\textsc{obj}.\textsc{art}3.\textsc{sg}	kopihoo\\
    \glt `The men cut a \textit{kopihoo} tree ...' (\ref{chap:23}:\ref{ex:text17-3})
\z
When the topical NP functions as the object, the subject NP or Pronoun follows the VC, see (\ref{ex:sketch-10}).
\ea\label{ex:sketch-10}
    O kaku vagana bona na pakupaku raara teo gohoho ge teo tarahi,\\
    \gll o	kaku	vagana	bona	na	paku{$\sim$}paku	ra=    ara	te=o	gohoho	ge	te=o	tarahi\\
    [\textsc{art}3.\textsc{sg}	kind	fish	\textsc{ana}.\textsc{sg}]\textsubscript{NP}	[\textsc{tm}2	\textsc{red}{$\sim$}do	1\textsc{pl}.\textsc{in}.\textsc{ipfv}=]\textsubscript{VC}    1\textsc{pl}.\textsc{in}.\textsc{pron}	\textsc{prep}=\textsc{art}3.\textsc{sg}	incoming.tide	or	\textsc{prep}=\textsc{art}3.\textsc{sg}	ebbing.tide\\
    \glt `This kind of fishing we do during the incoming tide or during the ebbing tide,' (\ref{chap:30}:\ref{ex:text23-3})
\z
Ditransitive clauses are frequent in this text collection, because in Teop the primary and secondary object NPs do not only refer to the recipient and the theme of a transfer,\footnote{Cf. \textcite[1--2]{Malchukov:2010wb}. For examples of \textit{hee} `give' see \textcite[487--500]{Mosel:2010bh} and \url{https://dictionaria.clld.org/units/teop-hee}.} but also to the patient and the instrument of an action or the product of an action and the material it is made of. The structure of ditransitive clauses is summarized in \tabref{tab:ditrans}.

\begin{table}
        \caption{Ditransitive clauses with a topic-comment structure.} \label{tab:ditrans}
\begin{tabularx}{\textwidth}{lllll}
        \lsptoprule
        \textsc{topic}&\textsc{vc}&\multicolumn{2}{c}{\textsc{non-topical arguments}}&\textsc{ref}\\
        \midrule
        subj: \textsc{h}-1\textsc{pl}.\textsc{in}.\textsc{pron}&VC\textsubscript{DTR}&1\textsuperscript{st} obj: \textsc{art} N&2\textsuperscript{nd} obj: \textsc{obj}.\textsc{art} N&(\ref{ex:sketch-11})\\
        subj: \textsc{h}-3\textsc{pl}.\textsc{pron}&VC\textsubscript{DTR}&1\textsuperscript{st} obj: \textsc{obj}.\textsc{art} N  &2\textsuperscript{nd} obj: \textsc{obj}.\textsc{art} N&(\ref{ex:sketch-12})\\
        1\textsuperscript{st} obj: \textsc{art} N&VC\textsubscript{DTR}&subj: 1\textsc{pl}.\textsc{in}.\textsc{pron}&2\textsuperscript{nd} obj: \textsc{obj}.\textsc{art} N&(\ref{ex:sketch-13})\\
        2\textsuperscript{nd}  obj: \textsc{art} N&VC\textsubscript{DTR}&subj: 3\textsc{pl}.\textsc{pron}&1\textsuperscript{st} obj: NP&(\ref{ex:sketch-14})\\
        \lspbottomrule
\end{tabularx}
\end{table}

The primary object NPs are treated like the object NP of transitive clauses. It is marked by the Basic Article, if the topical subject is formed by a 1\textsuperscript{st} or 2\textsuperscript{nd} Person Pronoun, see (\ref{ex:sketch-11}).
\ea\label{ex:sketch-11}
    ..., eara na tasu rara a koopu bona toraara.\\
    \gll e-ara	na	tasu	rara	a	koopu	bona	toraara\\
    \textsc{h}-1\textsc{pl}.\textsc{in}.\textsc{pron}	\textsc{tm}2	beat	1\textsc{pl}.\textsc{in}.\textsc{ipfv}	\textsc{art}3.\textsc{pl}	bamboo	\textsc{obj}.\textsc{art}2.\textsc{sg}	axe\\
    \glt `..., we beat the \textit{koopu} bamboos with an axe.' (\ref{chap:19}:\ref{ex:text14-21})
\z
But if the topical subject is formed by a 3\textsuperscript{rd} Person Pronoun or a NP, the primary object NP is marked by the Object Article, see (\ref{ex:sketch-12}).
\ea\label{ex:sketch-12}
    eori repaa kahu mana bono kara baneo bona kehaa\\
    \gll e-ori	re=  paa	kahu	mana	bono	kara	baneo	bona	kehaa\\
    \textsc{h}-3\textsc{pl}.\textsc{pron}	\textsc{consec}= [\textsc{tm}3	scrape	eventually]\textsubscript{VC}	\textsc{obj}.\textsc{art}3.\textsc{sg}	fibre	breadfruit\_tree	\textsc{obj}.\textsc{art}3.\textsc{sg}	shell\\
    \glt `they eventually scrape the breadfruit tree fibres with a shell (that they call \textit{sui}).' (\ref{chap:8}:\ref{ex:text4-20})
\z
The Basic Article only marks non-topical secondary object NPs, if both the subject and the primary object are formed by a 1\textsuperscript{st} or 2\textsuperscript{nd}  Person Pronoun.\footnote{There is no example in this text collection where both the subject and the primary object are formed by a 1\textsuperscript{st} or 2\textsuperscript{nd} person Pronoun, but see \textit{Ean toro hee anaa a moon.} `You must give me the girl.' in \textcite[490]{Mosel:2010bh}.}

When the primary object functions as the topic, the subject is formed by a Pronoun that directly follows the VC and the secondary object NP is marked by the Object Article irrespective of the person of the subject Pronoun, see (\ref{ex:sketch-13}).
\ea\label{ex:sketch-13}
    O pusu hiovu na paku raara bene ta boon.\\
    \gll o	pusu\_hiovu	na	paku	ra=  ara	bene	ta	boon\\
    \textsc{art}3.\textsc{sg}	husker	[\textsc{tm}2	make	1\textsc{pl}.\textsc{in}.\textsc{ipfv}=]\textsubscript{VC} 1\textsc{pl}.\textsc{in}.\textsc{pron}	\textsc{obj}.\textsc{art}1.\textsc{sg}	piece	mangrove\\
    \glt `We make the husker from a piece of mangrove wood.' (lit. `The husker make we') (\ref{chap:5}:\ref{ex:text01-52})
\z
The same clause structure is found with secondary object NPs functioning as the topic, see (\ref{ex:sketch-14}):
\ea\label{ex:sketch-14}
    O pihae vai o pereperee na paku riori bona bero a taba.\\
    \gll o	pihae	vai	o	pereperee	na	paku	ri=ori	bona	bero\footnotemark{}	a	taba\\
    \textsc{art}3.\textsc{sg}	coconut.frond	\textsc{dem}6	\textsc{art}3.\textsc{sg}	green	\textsc{tm}2	make	3\textsc{pl}.\textsc{ipfv}=3\textsc{pl}.\textsc{pron}	\textsc{obj}.\textsc{art}2.\textsc{sg}	many	\textsc{art}2.\textsc{sg}	thing\\
    \glt `The green coconut frond is used for many things.' (lit. `... make they many things.') (\ref{chap:5}:\ref{ex:text01-100})
\z\footnotetext{The word \textit{bero} `many' heads a NumPs in (\ref{ex:sketch-14}), but it also functions as the head of a VC, see (\ref{chap:13}:\ref{ex:text9-8}), the head of an AdjP, see (\ref{chap:26}:\ref{ex:text19-16}), and the head of a NP, see (\ref{chap:16}:\ref{ex:text11-12}).}
Thetic verbal clauses lack a topical argument as they introduce a new participant into the discourse or focus on the whole situation. Consequently, the arguments follow their VC, see (\ref{ex:sketch-15}).
\ea\label{ex:sketch-15}
    Evehee na teitei batana a sii vapara nae.\\
    \gll evehee	na	tei{$\sim$}tei	bata=na	a	sii	vapara	nae\\
    but	\textsc{tm}2	\textsc{red}{$\sim$}be.there	\textsc{cont}=3\textsc{sg}.\textsc{ipfv}	\textsc{art}2.\textsc{sg}	\textsc{dim}	partner	3\textsc{sg}.\textsc{poss}.\textsc{pron}\\
    \glt `But usually there is its partner.' (\ref{chap:35}:\ref{ex:text28-12})
\z
Clauses may lack subject and object NPs or Pronouns, if the context provides sufficient information to identify who or what is talked about, see \tabref{tab:zero}.
\section{The syntactic functions of phrases}\label{sec:synfunc}
The description of the syntactic functions of phrases is divided into two sections: \sectref{sec:phrasefuncclause} describes their functions as immediate clause constituents and \sectref{sec:subphrase} describes their functions as subordinate constituents of other phrases. Clause external topics are formed by NPs (24) and PrepPs (2).
\subsection{Phrases functioning as immediate clause constituents}\label{sec:phrasefuncclause}
Seven types of phrase function as immediate constituents of clauses, see \tabref{tab:phrasefunc}.


\begin{table}
        \caption{Phrases functioning as immediate clause constituents.} \label{tab:phrasefunc}
\begin{tabularx}{\textwidth}{lXlXl}
        \lsptoprule
        \textsc{phrase}&\multicolumn{4}{c}{\textsc{function as immediate clause constituents}}\\
        \cmidrule{2-5}
        &\textsc{predicate}&\textsc{subject/object}&\textsc{ob}\textsc{lique}&\textsc{adjunct}\\
        \midrule
        VC&(\ref{ex:sketch-1}), (\ref{ex:sketch-8})ff&-&-&-\\
        NP&(\ref{ex:sketch-2})&(\ref{ex:sketch-1}), (\ref{ex:sketch-3}), (\ref{ex:sketch-4}), (\ref{ex:sketch-5}), (\ref{ex:sketch-8})ff&-&-\\
        AdjP&(\ref{ex:sketch-4}), (\ref{ex:sketch-5})&-&-&-\\
        LocP&(\ref{ex:sketch-2})&-&(\ref{ex:sketch-16})&(\ref{ex:sketch-1})\\
        PrepP&(\ref{ex:sketch-6})&-&(\ref{ex:sketch-17})&(\ref{ex:sketch-10})\\
        anaphoric NumP&-&(\ref{ex:sketch-2}), (\ref{ex:sketch-6}), (\ref{chap:6}:\ref{ex:text2-35})&-&-\\
        anaphoric SubstP&-&(\ref{ex:sketch-2}), (\ref{ex:sketch-6}), (\ref{chap:29}:\ref{ex:text22-72})&-&-\\
        \lspbottomrule
\end{tabularx}
\end{table}

VCs only function as predicates. They are typically, but not obligatorily marked by particles that signify temporal relations, modality  and aspect, see \sectref{sec:tmmarker} and \sectref{sec:aspekt}. All other phrase types have more than one function. For the predicative function of NPs, AdjPs, LocPs and PrepPs see \sectref{sec:nonverbal}.

NPs function as topical and non-topical subjects of Non-Verbal Clauses, see \sectref{sec:nonverbal}, and in Verbal Clauses as topical and non-topical subjects and objects, see \sectref{sec:verbal}.

Anaphoric NumPs and SubstPs function as topical subjects and non-topial objects.

LocPs and PrepPs function as predicates and as oblique objects that are governed by a VC headed by \textit{tei} `be somewhere' or by a VC with the incorporated Preposition \textit{vo}\textsuperscript{1} `to' or \textit{vo}\textsuperscript{2} `like'.

Oblique objects formed by LocPs or PrepPs typically depend on VCs headed by \textit{tei} `be somewhere', see (\ref{ex:sketch-16}).

Oblique objects are governed by a VC headed by \textit{tei} `be somewhere' and by a VC with the incorporated preposition \textit{vo}\textsuperscript{1} `to' or \textit{vo}\textsuperscript{2} `like', and are expressed by an Adverb, LocP, or a PrepP. The LocPs or PrepPs governed by \textit{tei} `be somewhere' refer to any kind of location, be it inside a coconut (\ref{chap:5}:\ref{ex:text01-7}) or in the jungle, see (\ref{ex:sketch-16}).
\ea\label{ex:sketch-16}
    O kara kuskus na tei nana rasuu,\\
    \gll o	kara\_kuskus	na	tei	nana	rasuu\\
    \textsc{art}3.\textsc{sg}	passionflower	[\textsc{tm}2	be.there	3\textsc{sg}.\textsc{ipfv}]\textsubscript{VC}	jungle\\
    \glt `The passionflower is found in the jungle.' (\ref{chap:11}:\ref{ex:text7-39})
\z
VCs marked by the incorporated preposition \textit{vo} \textsc{goal} are mostly headed by a Verb denoting a kind of motion, see (\ref{ex:sketch-17}).
\ea\label{ex:sketch-17}
    ore paa kuu gunaha vo maa teo kasuana.\\
    \gll o=re	paa	kuu	gunaha	vo	maa	te=o	kasuana\\
    3\textsc{pl}.\textsc{pron}=\textsc{consec}	[\textsc{tm}3	fall	go.down	\textsc{goal}	hither]\textsubscript{VC}	\textsc{prep}=\textsc{art}3.\textsc{sg}	ground\\
    \glt `(When the coconut has become dry high up on a standing coconut palm,) then it falls down to the ground.' (\ref{chap:5}:\ref{ex:text01-3})
\z
The incorporated preposition \textit{vo}\textsuperscript{2} `like' is most frequently found with the cataphoric adverb \textit{en} `this' (\textsc{dem}3) and the anaphoric PrepP \textit{tee}=\textit{bona} `this' (\textsc{prep}={\allowbreak}\textsc{ana}.\textsc{sg}). Typical examples are (\ref{chap:5}:\ref{ex:text01-23}) and (\ref{chap:5}:\ref{ex:text01-45}).
\subsection{Subordinate phrases}\label{sec:subphrase}
VCs and the anaphoric NumPs and SubstPs are not found in subordinate functions, but all other phrases occur as subordinate phrases that are either embedded in the superordinate phrase or juxtaposed to it. \tabref{tab:subphrasefunc} below shows the distribution of six phrase types across six subordinate functions. A special phrase type is the Linker Phrase (LnkP) which only functions as a subordinate constituent. The subordinate functions of phrases are:
\begin{itemize}
    \item the core of  a VC, see (\ref{ex:sketch-34}) in \sectref{sec:verbcomplintro};
    \item the juxtaposed modifier of a Verb or Noun;
    \item the attribute of a NP;
    \item the inalienable possessor NP, see \sectref{sec:posscon};
    \item the complement of a Preposition in a PrepP;
    \item the complement of the Linker of a LnkP
\end{itemize}


\begin{table}
        \caption{Subordinate Phrases.} \label{tab:subphrasefunc}
\small
\begin{tabularx}{\textwidth}{l@{~~}p{9mm}Q@{~~}lQ@{}Qp{15mm}@{}}
        \lsptoprule
        \textsc{phrase}&\multicolumn{6}{c}{\textsc{functions as subordinate constituents}}\\
        \cmidrule{2-7}
        &\textsc{core of VC}&\textsc{modifier of V or N}&\textsc{attr. of NP}&\textsc{inalien. possessor}&\textsc{compl. of \mbox{preposition}}&\textsc{compl. of linker}\\
        \midrule
        NP&(\ref{ex:sketch-18})&-&(\ref{ex:sketch-25})&(\ref{ex:sketch-23}), (\ref{ex:sketch-25})&\mbox{(\ref{ex:sketch-19})-(\ref{ex:sketch-22}),~(\ref{ex:sketch-24})}&-\\
        AdjP&-&-&(\ref{ex:sketch-3}), (\ref{ex:sketch-14}), (\ref{ex:sketch-26})&-&-&-\\
        NumP&-&-&(\ref{ex:sketch-2})&-&-&-\\
        LocP&-&-&-&-&-&(\ref{ex:sketch-23})\\
        PrepP&(\ref{ex:sketch-19})&-&(\ref{ex:sketch-22}), (\ref{ex:sketch-24})&-&-&(\ref{ex:sketch-20})-(\ref{ex:sketch-22})\\
        LnkP&(\ref{ex:sketch-20})&(\ref{ex:sketch-21})-(\ref{ex:sketch-23})&-&-&-&-\\
        \lspbottomrule
\end{tabularx}
\end{table}

NPs and PrepP that function as the predicate of a non-verbal clause can be negated by the Double Negator \textit{sa(ka)} ... \textit{haa} (\textsc{neg}1 ... \textsc{neg}2) and then function as the core of a VC. In a text about poisoning fish with roots, the author emphasizes that they use many roots, see (\ref{ex:sketch-18}).
\ea\label{ex:sketch-18}
Saka o peho ana haana to tasu raara,\\
\gll saka	o	peho	ana	haa=na	to	tasu	ra=ara\\
     [\textsc{neg}1	\textsc{art}3.\textsc{sg}	one3.\textsc{sg}	root	\textsc{neg}=3\textsc{sg}.\textsc{ipfv}]\textsubscript{VC}	\textsc{rel}	beat	1\textsc{pl}.\textsc{in}.\textsc{ipfv}=1\textsc{pl}.\textsc{in}.\textsc{pron}\\
\glt `It is not a single root that we beat (and throw into the pool), ...' (\ref{chap:37}:\ref{ex:text30-15})
\z
Similarly, a PrepP functions as the core of a negated VC to emphasize that turtle netting is not done at the edge of the reef, see (\ref{ex:sketch-19}).
\ea\label{ex:sketch-19}
ei saka tea irihi haana,\\
\gll ei	saka	te=a	irihi	haa=na\\
     \textsc{dem}2	[\textsc{neg}1	[\textsc{prep}=\textsc{art}2.\textsc{sg}	edge.of.reef]\textsubscript{PrepP}	\textsc{neg}2=3\textsc{sg}.\textsc{ipfv}]\textsubscript{VC}\\
\glt `(We know the place where the net is thrown or where the turtle netting is done,) this (is) not at the edge of the reef, ...' (\ref{chap:28}:\ref{ex:text21-11})
\z
Thirdly, a LnkP with a PrepP as its complement holds the position of the core of an affirmative VC and means `be like sth', see (\ref{ex:sketch-20}).
\ea\label{ex:sketch-20}
A kokoroto na vaa tea ahio nana,\\
\gll a kokoroto na vaa te=a ahio nana\\
    \textsc{art}2.\textsc{sg} kokoroto \textsc{tm}2 \textsc{lnk} \textsc{prep}=\textsc{art}2.\textsc{sg}	ahio 3\textsc{sg}.\textsc{ipfv}\\
    \glt `The \textit{kokoroto} is similar to the \textit{ahio} lizard (but the \textit{ahio} is green, while the \textit{kokoroto} is brown.)' (See \url{https://dictionaria.clld.org/units/teop-kokoroto})
\z
In this text collection the LnkP is used as a modifier of a Verb, see (\ref{ex:sketch-21}),
\ea\label{ex:sketch-21}
o sii vahara vasu beori ravun vaa tea vaakororovi rori.\\
\gll o	sii	vahara	vasu	be=ori	ravun	vaa	te=a	vaakororovi	rori\\
     \textsc{art}2.\textsc{pl}	\textsc{dim}	small.\textsc{pl}	stone	when=3\textsc{pl}.\textsc{pron}	be.heaped.up	\textsc{lnk}	\textsc{prep}=\textsc{art}2.\textsc{sg}	circle	3\textsc{pl}.\textsc{ipfv}\\
\glt `(When we are about to catch wrasses, we search for) little stones being heaped up in a circle.' (\ref{chap:35}:\ref{ex:text28-6})
\z
or as a modifier of a Noun\footnote{In (\ref{chap:5}:\ref{ex:text01-180}) \textit{mene} \textit{vaanoto} \textit{suraa} \textit{vaa} \textit{tea} \textit{taotao} `in order to light the fire for baking (food) in the earth oven' the LnkP modifies the Noun \textit{suraa} `fire' which functions as the complement of the transitive verb \textit{vaanoto} `light'.}, see (\ref{ex:sketch-22}).
\ea\label{ex:sketch-22}
    O paranda o peho taba vaa tea inu teara ...\\
    \gll o	paranda	o	peho	taba	vaa	te=a	inu	te=ara\\
    \textsc{art}3.\textsc{sg}	floor	\textsc{art}3.\textsc{sg}	\textsc{indef}3.\textsc{sg}	thing	\textsc{lnk}	\textsc{prep}=\textsc{art}2.\textsc{sg}	house	\textsc{prep}=1\textsc{pl}.\textsc{in}.\textsc{pron}\\
    \glt `The floor is a thing in our house, ...' (\ref{chap:20}:\ref{ex:text15-1})
\z
Similar to the PrepP in (\ref{ex:sketch-22}), LocPs are linked as a modifier to a Noun, see (\ref{ex:sketch-23}).
\ea\label{ex:sketch-23}
    a vana vaa komana overe\\
    \gll a	vana	vaa	koma-n=a	overe\\
    \textsc{art}3.\textsc{pl}	germinating.leaf	\textsc{lnk}	inside-3\textsc{sg}.\textsc{poss}=\textsc{art}2.\textsc{sg}	coconut\\
    \glt `the germinating leaves inside the coconut' (\ref{chap:5}:\ref{ex:text01-24})
\z
In NPs referring to the stages of work processes, the Linker connects a Complement Clause to the head, e.g. \textit{a kiu vaa tea} V... `the work of doing V', see \sectref{sec:depclaus}.
When a PrepP introduced by \textit{te} (\textsc{prep}) directly follows a NP as its attribute, it always refers to an alienable possessor, see (\ref{ex:sketch-24}).
\ea\label{ex:sketch-24}
    A kara bao, eve a tabaan teo burumakao.\\
    \gll a	kara	bao	eve	a	tabaan	te=o	burumakao\\
    \textsc{art}3.\textsc{pl}	string	grass	3\textsc{sg}.\textsc{pron}	[\textsc{art}2.\textsc{sg}	food	\textsc{prep}=\textsc{art}2.\textsc{pl}	cow]\textsubscript{NP}\\
    \glt `Grass is the food of cows.' (\ref{chap:11}:\ref{ex:text7-11})
\z
In the Inalienable Possessive Constructions (\ref{ex:sketch-23}) and (\ref{ex:sketch-25}), the head Noun is marked by the suffixed Possessive Marker (\textsc{poss}) which indexes the person and number category of the following possessor NP, see \sectref{sec:posscon}.
Attributive NPs modify the meaning of a preceding NP marked by \textit{vai} `this' (\textsc{dem}6), see (\ref{ex:sketch-25}).\footnote{Similar examples are (\ref{chap:6}:\ref{ex:text2-2}) and (\ref{chap:7}:\ref{ex:text3-4})}
\ea\label{ex:sketch-25}
    O naono vai o kopihoo na paku riori bono aga no sinivi.\\
    \gll o	naono	vai	o	kopihoo	na	paku	ri=ori	bono	aga	n=o	sinivi\\
    \textsc{art}3.\textsc{sg}	tree	\textsc{dem}6	\textsc{art}3.\textsc{sg}	kopihoo	\textsc{tm}2	make	3\textsc{pl}.\textsc{ipfv}=3\textsc{pl}.\textsc{pron}	\textsc{obj}.\textsc{art}3.\textsc{sg}	float	3\textsc{sg}.\textsc{poss}=\textsc{art}3.\textsc{sg}	canoe\\
    \glt `The \textit{kopihoo} tree is used for making the float of the outrigger canoe.' (\ref{chap:23}:\ref{ex:text17-2})
\z
NPs followed by an attributive AdjP are optionally marked by \textit{vai} `this' (\textsc{dem}6) to distinguish the attriutive AdjP from the predicative one, see (\ref{ex:sketch-26}), and compare it with (\ref{ex:sketch-4}) above.
\ea\label{ex:sketch-26}
    A koopu vai a mataa ae bebeahu na teitei nana rasuu.\\
    \gll a	koopu	vai	a	mataa	ae	bebeahu	na	tei{$\sim$}tei	nana	rasuu\\
    \textsc{art}3.\textsc{pl}	bamboo	\textsc{dem}6	\textsc{art}3.\textsc{pl}	good	and1	long	\textsc{tm}2	\textsc{red}{$\sim$}be.there	3\textsc{sg}.\textsc{ipfv}	jungle\\
    \glt `The good and long \textit{koopu} bamboos are in the jungle.' (\ref{chap:19}:\ref{ex:text14-6})
\z
\section{The classification of content words and Numerals}\label{sec:numerals}
\subsection{Introduction}
This section deals with Nouns (N), Verbs (V), Adjectives (Adj), Cardinal Numerals (Num)\footnote{Other Numerals are derived from Cardinal Numerals, e.g. \textit{vaa-buaku} `second' (\ref{chap:20}:\ref{ex:text15-18}), \textit{bubuaku} `in pairs', and \textit{va-buaku} `two together' (\ref{chap:29}:\ref{ex:text22-63}).} and Manner Adverbs (Adv) that are derived by the prefix \textit{va}- `\textsc{advr}'. Their classification is based on the combination of two criteria:
\begin{enumerate}
    \item their function as heads of VCs, NPs, AdjPs, and NumPs;
    \item their compatibility with directly following content words as their modifiers.
\end{enumerate}
LocPs are not considered in this section, because they are exclusively headed by Nouns referring to a specific place. A special subclass of Nouns which denote specific kinds of plurality are dealt with in \sectref{sec:classnphead}. The most frequent Noun of this kind is the Plural Marker \textit{maa} `\textsc{plm}' denoting distributional plurality. For details of the distinction of Teop Nouns, Verbs, and Adjectives see \textcite{Mosel:2017as}, for lexical flexibility in Oceanic languages \textcite{Lier:2016ka}, and for word classification in Austronesian languages \textcite{Mosel:2023sy}.
\subsection{The function of content words and Numerals as phrasal heads}\label{sec:funcnumhead}
Verbs and Adjectives occur as the head of VCs, NPs, and AdjPs, whereas Nouns only head VCs and NPs. But Verbs differ from Adjectives in their selection of modifiers, see \sectref{sec:modjux}. Cardinal Numerals function as the head of VCs and NumPs. Adverbs are not found in the position of a phrasal head, see \tabref{tab:numhead}.


\begin{table}
    \caption{Content words and Numerals functioning as the head of VCs, NPs, AdjP, and NumPs.} \label{tab:numhead}
\begin{tabularx}{\textwidth}{lXlXl}
    \lsptoprule
    \textsc{word classes}&VC&NP&AdjP&NumP\\
    \midrule
        V&(\ref{ex:sketch-27b})&(\ref{ex:sketch-27a}), (\ref{ex:sketch-28}), (\ref{ex:sketch-29})&(\ref{ex:sketch-30b}), (\ref{ex:sketch-38})&-\\
        N&(\ref{ex:sketch-31})&(\ref{ex:sketch-28}), (\ref{ex:sketch-27b})&-&-\\
        Adj&(\ref{ex:sketch-28})&(\ref{ex:sketch-81b})&(\ref{ex:sketch-30a}), (\ref{ex:sketch-81a})&-\\
        Num&(\ref{ex:sketch-32})&-&-&(\ref{ex:sketch-27a}), (\ref{ex:sketch-33})\\
    \lspbottomrule
\end{tabularx}
\end{table}

NPs headed by a Verb have the same syntactic properties as NPs headed by a Noun with respect to attributive NumPs, see (\ref{ex:sketch-27}).
\ea\label{ex:sketch-27}
\ea\label{ex:sketch-27a}
    A kukan bau gea dorana a hovu,\\
    \gll a	kukan	bau	ge=a	dorana	a	hovu\\
    [\textsc{art}2.\textsc{sg}	three	perhaps	or=\textsc{art}2.\textsc{sg}	four]\textsubscript{NumP}	[\textsc{art}2.\textsc{sg}	husk]\textsubscript{NP}\\
    \glt `Perhaps three or four husking movements, and then the coconut is completely husked.'
\ex\label{ex:sketch-27b}
    a overe repaa tahovu vakavara.\\
    \gll a	overe	re=paa	ta-hovu	va-kavara\\
    \textsc{art}2.\textsc{sg}	coconut	\textsc{consec}=\textsc{tm}3	\textsc{acaus}-husk	\textsc{advr}-complete\\
    \glt `(and) then the coconut is completely husked.' (\ref{chap:5}:\ref{ex:text01-62})
\z
\z
When the NP head is formed by a transitive Verb, it functions as the possessum of an Inalienable Possessive Construction (\sectref{sec:posscon}), see (\ref{ex:sketch-28}).
\ea\label{ex:sketch-28}
    A koro no komano sinivi potee na varuruva koana\\
    \gll a	koro	n=o	koma-n=o	sinivi	potee	na	varuruva	koa=na\\
    [\textsc{art}2.\textsc{sg}	carve.out	3\textsc{sg}.\textsc{poss}=\textsc{art}3.\textsc{sg}	inside-3\textsc{sg}.\textsc{poss}=\textsc{art}3.\textsc{sg}	canoe	potee]\textsubscript{NP}	[\textsc{tm}2	easy	just=3\textsc{sg}.\textsc{ipfv}]\textsubscript{VC}\\
    \glt `The carving of the inside of the \textit{potee} canoe is quite easy.' (\ref{chap:24}:\ref{ex:text18-14}) 
\z
But within the NP, a Verb is modified by an Adverb and not by an Adjective, see (\ref{ex:sketch-29}).
\ea\label{ex:sketch-29}
    ahiki ta popo vatamaka\\
    \gll ahiki ta popo va-tamaka\\
    not.exist \textsc{nspec}2.\textsc{sg} live \textsc{advr}-sad\\
    \glt `there was not any sad life' \\\hfill(See \url{https://dictionaria.clld.org/units/teop-vatamaka})
\z
Similarly, Verbs heading an AdjP are modified by adverbs as, for instance, \textit{kare} `recently', see (\ref{ex:sketch-30}).
\ea\label{ex:sketch-30}
\ea\label{ex:sketch-30a}
    O baneo vai o visoasi,\\
    \gll o	baneo	vai	o	visoasi\\
    \textsc{art}3.\textsc{sg}	breadfruit\_tree	\textsc{dem}6	\textsc{art}3.\textsc{sg}	young\\
    \glt `The young breadfruit tree'
\ex\label{ex:sketch-30b}
    o baneo vai o kare pura,\\
    \gll o	baneo	vai	o	kare	pura\\
    \textsc{art}3.\textsc{sg}	breadfruit\_tree	\textsc{dem}6	[\textsc{art}3.\textsc{sg}	recently	grow]\textsubscript{AdjP}\\
    \glt `is a breadfruit tree that has only recently started growing, ...' (\ref{chap:8}:\ref{ex:text4-1})
\z
\z
Nouns heading a VC are modified by Adjectives, see (\ref{ex:sketch-31}).
\ea\label{ex:sketch-31}
O naono vai o sinivi na taba beera kurusu kiri nana bona tavaan,\\
\gll o	naono	vai	o	sinivi	na	taba	beera	kurusu	ki=ri	nana	bona	tavaan\\
     \textsc{art}3.\textsc{sg}	tree	\textsc{dem}6	\textsc{art}3.\textsc{sg}	sinivi	[\textsc{tm}2	thing	important	very	\textsc{dat}=3\textsc{pl}.\textsc{objm}	3\textsc{sg}.\textsc{ipfv}]\textsubscript{VC}	\textsc{obj}.\textsc{art}2.\textsc{sg}	people\\
\glt `This canoe tree is a very important thing for the people, ...' (\ref{chap:7}:\ref{ex:text3-4})
\z
Adjectives function as the head of AdjPs (\ref{ex:sketch-30a}), VCs (\ref{ex:sketch-28}), and Inalienable Possessive NPs, see \sectref{sec:posscon}.

Cardinal numerals function as the head of VCs, see (\ref{ex:sketch-32}),
\ea\label{ex:sketch-32}
    Enam paa dorana vaevuru.\\
    \gll e-nam paa dorana vaevuru\\
    \textsc{h}-1\textsc{pl}.\textsc{ex}.\textsc{pron} \textsc{tm}3 four already\\
    \glt `We were already four.' (See \url{https://dictionaria.clld.org/units/teop-dorana})
\z
and as the head of NumPs functioning as an attribute or an anaphoric argument, see (\ref{ex:sketch-33}).
\ea\label{ex:sketch-33}
\ea\label{ex:sketch-33a}
    Teo dorana o pakapaka,\\
    \gll te=o dorana o pakapaka\\
    \textsc{prep}=\textsc{art}3.\textsc{sg} four \textsc{art}3.\textsc{sg}	flipper\\
    \glt `As for the four flippers (of the turtle),'
\ex\label{ex:sketch-33b}
    o buaku vaa kahoo nae na dao raara bono huahua.\\
    \gll o buaku vaa kahoo nae na dao ra=ara bono hua{$\sim$}hua\\
    \textsc{art}3.\textsc{sg} two \textsc{lnk} head-	3\textsc{sg}.\textsc{poss}.\textsc{pron} \textsc{tm}2 call 1\textsc{pl}.\textsc{in}.\textsc{ipfv}=1\textsc{pl}.\textsc{in}.\textsc{pron} \textsc{obj}.\textsc{art}3.\textsc{sg} \textsc{red}{$\sim$}paddle\\
    \glt `the two near its head, we call paddles.' (\ref{chap:26}:\ref{ex:text19-6})
\z
\z
In (\ref{ex:sketch-33b}) the NumP \textit{o buaku} `the two' is classified as an anaphoric NumP functioning as a topical primary object, because the selection of its Article \textit{o} `\textsc{art}3.\textsc{sg}' is determined by the head of the preceding NP \textit{o dorana o pakapaka} `the four flippers'.
\subsection{The modification of content words by other juxtaposed content words}\label{sec:modjux}
The term modifier is used for content words that directly follow another content word and characterize the concept denoted by the preceding word. Not considered as modifiers are content words functioning as a secondary predication, e.g. \textit{avuhu mataa} `smell good' (\ref{chap:5}:\ref{ex:text01-139}), or as a juxtaposed complement of a relational Noun, e.g. NN: \textit{kaku vagana} `kind of fishing' (\ref{ex:sketch-10}), or a transitive Verb, e.g. VN: \textit{paku sinivi} `make (a) canoe' (\ref{chap:7}:\ref{ex:text3-14}), VV: \textit{vaanagu koava} `resist rotting' (\ref{chap:16}:\ref{ex:text11-8}), VNum: \textit{gae kukan} `lift three (bamboo strips)' (\ref{chap:19}:\ref{ex:text14-35}).
Verbs, Nouns, Adjectives and numerals are distinct from one another with respect to their modification by content words. While Numerals do not have any modifiers, Verbs, Nouns and Adjectives differ in their selection of modifying content words, see \tabref{tab:juxtaposed}.


\begin{table}
    \caption{The compatibility of Verbs, Nouns and Adjectives with juxtaposed modifiers.} \label{tab:juxtaposed}
\begin{tabularx}{.8\textwidth}{XXXXl}
    \lsptoprule
    &\multicolumn{4}{c}{\textsc{Modifiers}}\\
    \cmidrule{2-5}
    &V&N&Adj&Adv\\
    \midrule
        V&(\ref{chap:5}:\ref{ex:text01-3})&(\ref{chap:22}:\ref{ex:text16-23})&-&(\ref{chap:10}:\ref{ex:text6-30})\\
        N&(\ref{chap:16}:\ref{ex:text11-1})&(\ref{chap:5}:\ref{ex:text01-20})&(\ref{chap:7}:\ref{ex:text3-4})&-\\
        Adj&(\ref{chap:30}:\ref{ex:text23-18})&-&-&(\ref{chap:27}:\ref{ex:text20-14})\\
    \lspbottomrule
\end{tabularx}
\end{table}

Verbs can be modified by Verbs, Nouns and Adverbs, see \tabref{tab:modverbs}.

\begin{table}
    \caption{The modification of Verbs.} \label{tab:modverbs}
\begin{tabularx}{\textwidth}{Xllll}
    \lsptoprule
    &\textsc{text}&\textsc{gloss}&\textsc{translation}&\textsc{ref}\\
    \midrule
        V V&\textit{kuu gunaha}&fall go.down&`fall down'&(\ref{ex:sketch-17})\\
        V N&\textit{paku vaagum}&do group&`do sth in a group'&(\ref{chap:22}:\ref{ex:text16-23})\\
        V Adv&\textit{pura va-mataa}&grow \textsc{advr}-good&`grow well'&(\ref{chap:10}:\ref{ex:text6-30})\\
    \lspbottomrule
\end{tabularx}
\end{table}

Nouns only differ from Verbs in being modified by an Adjective instead of an Adverb, see \tabref{tab:modnouns}.

\begin{table}
    \caption{The modification of Nouns.} \label{tab:modnouns}
\begin{tabularx}{\textwidth}{Xllll}
    \lsptoprule
    &\textsc{text}&\textsc{gloss}&\textsc{translation}&\textsc{ref}\\
    \midrule
        N V&\textit{inu goro{$\sim$}goroho}&house \textsc{red}{$\sim$}sleep&`sleeping house'&(\ref{chap:16}:\ref{ex:text11-1})\\
        N N&\textit{paka overe }&leaf coconut.palm&`coconut leaf'&(\ref{chap:5}:\ref{ex:text01-20})\\
        N Adj&\textit{taba beera}&thing big/important&`important thing'&(\ref{ex:sketch-31})\\
    \lspbottomrule
\end{tabularx}
\end{table}

Similar to Verbs, Adjectives are modified by Verbs and Adverbs, but it is only the three verbs \textit{oha} `exceed,' \textit{potee} `be like,' and \textit{vaapeha} `be the same' and a few adverbs. In contrast to Verbs, AdjPs are not modified by Nouns, see \tabref{tab:adjmod}.


\begin{table}[t]
    \caption{Adjectives modified by Verbs and Adverbs.} \label{tab:adjmod}
\begin{tabularx}{\textwidth}{Xllll}
    \lsptoprule
    &\textsc{text}&\textsc{gloss}&\textsc{translation}&\textsc{ref}\\
    \midrule
        \multirow{2}{*}{Adj V}&\textit{beera oha}&big exceed&`bigger than'&(\ref{chap:30}:\ref{ex:text23-18})\\
        &\textit{beera potee}&big be.like&`as big as'&(\ref{chap:24}:\ref{ex:text18-5})\\\midrule
        \multirow{2}{*}{Adj Adv}&\textit{bebeahu va-antee}&long \textsc{advr}-enough&`long enough'&(\ref{chap:27}:\ref{ex:text20-14})\\
        &\textit{paru va-pereperee}&black \textsc{advr}-green&`greenish black'&\textsc{mted}\footnote{\url{https://dictionaria.clld.org/units/teop-paru_vapereperee} and \url{https://dictionaria.clld.org/units/teop-vapereperee}.}\\
    \lspbottomrule
\end{tabularx}
\end{table}

The selection of modifiers is independent of the syntactic function of the modified content word. In (\ref{ex:sketch-29}) the Verb \textit{popo} `live' heads a NP and is modified by the adverb \textit{vatamaka} `sadly', and in (\ref{ex:sketch-31}) the prototypical Noun \textit{taba} `thing' functions as the head of a VC and is modified by the Adjective \textit{beera} `big, important'.
\section{The morphology}\label{sec:morphology}
The inflectional morphology is restricted to the Articles (\sectref{sec:paraarticles}), the Numerals \textit{peha}\textsuperscript{1}/\textit{peho}\textsuperscript{1} `one' and bua/buo `two', the Indefinite Determiner \textit{peha}\textsuperscript{2}/\textit{peho}\textsuperscript{2}, the Determiner \textit{meha}/\textit{meho} `other, some' within the NP (\sectref{sec:preheadnp}), to the person forms of the Object Marker (\sectref{sec:verbcomplintro}) and the Imperfective Aspect Marker in the VC (\sectref{sec:aspekt}), and to the Pronouns (\sectref{sec:pronouns}) and the Possessive Marker in Inalienable Possessive Constructions (\sectref{sec:posscon}).

The derivational morphology is restricted to seven affixes and reduplication:\footnote{For examples see \textcite[258--260]{Mosel:2007ng} or  Search (Words extra) in the \textsc{mted} for the Gloss. If you search for causatives, insert ``{\^{}}CAUS'', the graph ``{\^{}}'' means: the beginning of the entry.}
\begin{itemize}
    \item \textsc{acaus}-: the Anticausative prefix \textit{ta}- derives intransitive Verbs from transitive ones.
    \item -\textsc{adjr}: the suffix -\textit{ana} derives Adjectives from reduplicated Nouns. The form \textsc{red}{$\sim$}N\textit{-ana} means `full of N', e.g., \textit{usu{$\sim$}usu-ana} `full of lice'.
    \item \textsc{advr}-: the Prefix \textit{va}- derives Adverbs from Adjectives and some stative Verbs
    \item \textsc{caus}-: the Prefix \textit{vaa}-, \textit{va}- derives transitive causative Verbs and some intransitive Verbs with various meanings.
    \item -\textsc{derel}: the Derelational suffix -\textit{na} derives non-relational nouns from relational nouns, see \sectref{sec:posscon}.
    \item \textsc{mult}-: the Prefix \textit{va}- derives Verbs denoting multiple actions, including reciprocal actions.
    \item \textsc{ord}-: the Ordinal prefix \textit{vaa}- derives Ordinal Numerals from Cardinal Numerals.
    \item \textsc{red}{$\sim$}: Reduplication signifies various kinds of plurality and derives Adjectives from Nouns or other adjectives, e.g. \textit{pere{$\sim$}peree} `green' from \textit{peree} `unripe'.
\end{itemize}
In accordance with the multifunctionality of Verbs and Nouns, Teop lacks morphemes that derive Nouns from Verbs or Verbs from Nouns.
\section{The Verb Complex (VC)}\label{sec:vc}
\subsection{Introduction}\label{sec:verbcomplintro}
The VC is the only type of phrase that exclusively functions as a predicate, see \sectref{sec:phrasefuncclause}. The head of canonical VCs is formed by a Verb, a Noun or an Adjective; for non-canonical VCs see \sectref{sec:subphrase}. The sequential order of functional and lexical words in VCs is fixed, see (\ref{ex:sketch-34}).
\ea\label{ex:sketch-34}
    \textup{\textsc{neg}1 \textsc{tm}	[Adv\textsubscript{1} HEAD X Adv\textsubscript{2}]\textsubscript{\textsc{core}} 	\textsc{neg2} \textsc{appl} Prep \textsc{objm} Adv\textsubscript{3} \textsc{ipfv}/\textsc{pfv}}\\
\z
In (\ref{ex:sketch-34}) X stands for content words that follow the head as modifiers, complements or secondary predications, see \sectref{sec:modjux} and the modifying LnkP, see \sectref{sec:subphrase}. The abbreviations Adv\textsubscript{1}, Adv\textsubscript{2}, and Adv\textsubscript{3} indicate the three positions of Adverbs. The first position is held by \textit{pahin} `immediately' (\ref{chap:22}:\ref{ex:text16-24}) or \textit{manin} `ever' in negated VCs, see (\ref{chap:11}:\ref{ex:text7-63}).\footnote{Further examples of Adverbs in this position are \textit{pate}\textsuperscript{2} `together' and \textit{van} `together', search for these Headwords in the MTED.}
In contrast, the second position can be held by more than one adverb, see (\ref{ex:sketch-35}).
\ea\label{ex:sketch-35}
O koverau saka tau pura vatotobin tamuana bata haana.\\
\gll o	koverau	saka	tau	pura	va-totobin	tamuana	bata	haa=na\\
     \textsc{art}3.\textsc{sg}	koverau	\textsc{neg}1	\textsc{tm}5	grow	\textsc{advr}-straight	always	\textsc{cont}	\textsc{neg}2=3\textsc{sg}.\textsc{ipfv}\\
\glt `The \textit{koverau} bamboo might not always continuously grow straight.' (\ref{chap:10}:\ref{ex:text6-2})
\z
The third Adv-position can be held by \textit{maa} `hither' or \textit{nao} `thither' and roho `before', which may occur together as in (\ref{ex:sketch-37}) \textit{deedee} \textit{maa} \textit{roho} `brought be here'.

The functional words of the VC are:
\begin{enumerate}
    \item the Double Negator (\ref{ex:sketch-18}), (\ref{ex:sketch-19}), (\ref{ex:sketch-35});
    \item the pre-head TM-Markers signifying temporal relations and modality, e.g. \textit{tau} \textsc{tm}5 (\ref{ex:sketch-35}), see \sectref{sec:tmmarker};
    \item the post-head Applicative Particle \textit{ni} \textsc{appl} signifying that the VC is transitive, see \sectref{sec:applart}; 
    \item post-head incorporated Prepositions which increase the valency of the VC core, see \sectref{sec:incorporated};
    \item the post-head Object Markers (\textsc{objm}) which index direct and primary object NPs and Pronouns, i.e. the  uninflected Object Marker \textit{a} \textsc{objm} (\ref{ex:sketch-47}), \textit{ara} 1\textsc{pl}.\textsc{in}.\textsc{objm} (\ref{ex:sketch-43}), and \textit{ri} 3\textsc{pl}.\textsc{objm} (\ref{ex:sketch-31}), (\ref{ex:sketch-39c}), (\ref{ex:sketch-68}), (\ref{ex:sketch-70});
    \item the post-head Aspect Markers (\textsc{ipfv}, \textsc{pfv}) occur in the last position of the VC, see \sectref{sec:aspekt}.
\end{enumerate}
There are three kinds of Hybrid VCs:
\begin{enumerate}
    \item a negated VC in which the core is formed by a NP or PrepPs, see \sectref{tab:subphrasefunc};
    \item an affirmative VC with a LnkP as its core, see \sectref{tab:subphrasefunc}, and 
    \item a VC in which the core is formed by Noun and the Demonstrative \textit{vai} `this' (\textsc{dem6}), and then beyond its borderline followed by an modifying AdjP or Relative Clause. Examples are found in (\ref{chap:16}:\ref{ex:text11-11}), (\ref{chap:20}:\ref{ex:text15-8}) and (\ref{chap:31}:\ref{ex:text24-5}).
\end{enumerate}
In certain contexts such as Imperative Constructions, Adverbial Clauses, and Relative Clauses the VC may only consist of a single Verb, see (\ref{chap:26}:\ref{ex:text19-40}), (\ref{chap:5}:\ref{ex:text01-4}), and (\ref{ex:sketch-75}).

\subsection{The Applicative Particle}\label{sec:applart}
Irrespective of the valency of the VC core, the Applicative Particle marks the valency of a VC as transitive, see \textcite[493--496]{Mosel:2010bh}\footnote{Search (Words) in the \textsc{mted} for the \textsc{mwe}  ``vi - appl''  and find 123 entries with the Part of Speech ``vt.constr''. Then search for the \textsc{mwe} ``vd - appl'' and find 14 entries with the Part of Speech ``vt.constr''.}, see (\ref{ex:sketch-36}).

\ea\label{ex:sketch-36}
A ba suburaara na booboha overe ni roho bona vasu.\\
\gll a	ba	subu-raara	na	booboha	overe	ni	roho	bona	vasu\\
     \textsc{art}2.\textsc{sg}	\textsc{kin}.\textsc{plm}	grandparent-1\textsc{pl}.\textsc{in}.\textsc{poss}.\textsc{pron}	[\textsc{tm}2	[[break\textsubscript{\textsc{tr}}	coconut]\textsubscript{\textsc{intr}}	\textsc{appl}]\textsubscript{\textsc{tr}}	before]\textsubscript{VC}	[\textsc{obj}.\textsc{art}2.\textsc{sg}	stone]\textsubscript{\textsc{ins}}\\
    \glt `Our grandparents broke coconuts (by hitting them) with a stone, ...' (\ref{chap:5}:\ref{ex:text01-69})
\z
The Applicative Particle marks the valency of a VC as transitive. In (\ref{ex:sketch-36}) the transitive Verb \textit{booboha} `break' is followed by \textit{overe} `coconut' functioning as its complement. This intransitive construction \textit{booboha} \textit{overe} `break coconuts' is then transitivized by the Applicative Particle \textit{ni} \textsc{appl}, e.g. \textit{booboha} \textit{overe} \textit{ni} `break coconuts with sth.' Further examples of this frequently used construction are listed in \tabref{tab:transitivized}.

\begin{table}
    \caption{Transitive VCs of the structure V N\textsubscript{\textsc{patient}} \textsc{appl} NP\textsubscript{\textsc{instrument}}.} \label{tab:transitivized}
\begin{tabularx}{\textwidth}{Xll}
    \lsptoprule
    \textsc{text}&\textsc{translation}&\textsc{ref}\\
    \midrule
        \textit{kuukuhi} \textit{overe} \textit{ni}&`grate coconuts with (a shell)'&(\ref{ex:sketch-1})\\
        \textit{vaanoto} \textit{suraa} \textit{ni}&`light fire with (the fibrous husk)'&(\ref{chap:5}:\ref{ex:text01-119})\\
        \textit{vigu} \textit{inu} \textit{ni}&`use (the sago palm) for building houses'&(\ref{chap:6}:\ref{ex:text2-54})\\
        \textit{isuisuvu} \textit{ruene} \textit{ni}&`use (the \textit{kururu} bamboo) for fetching water'&(\ref{chap:10}:\ref{ex:text6-24})\\
        \textit{paku} \textit{vavaatoa} \textit{ni}&`do kite fishing with (\textit{ariono} vines)'&(\ref{chap:11}:\ref{ex:text7-9})\\
    \lspbottomrule
\end{tabularx}
\end{table}

If a VC is headed by the ditransitive Verb like \textit{hee} `give sb. sth.' and marked by \textit{ni}, the VC is transitive, e.g. \textit{hee} \textit{hiava} \textit{ni} [give go.up \textsc{appl}] `reach sth. up', (\ref{chap:6}:\ref{ex:text2-40}), (\ref{chap:17}:\ref{ex:text12-37}).
\subsection{Incorporated Prepositions}\label{sec:incorporated}
The incorporation of Prepositions in VCs is a means of increasing the valency of the VC core (see \cite{Mosel:2000pm}). The VC is mostly headed by a Verb, see \tabref{tab:vcverbincorp}.

\begin{table}
    \caption{VCs headed by a Verb with incorporated prepositions.} \label{tab:vcverbincorp}
\fittable{
\begin{tabular}{llll}
    \lsptoprule
    \textsc{text}&\textsc{gloss of prep}&\textsc{translation}&\textsc{ref}\\
    \midrule
        \textit{neba} \textit{ki}&\textsc{dat}&`plant sth. for (cows)'&(\ref{ex:sketch-39c})\\
        \textit{sue} \textit{ki}&\textsc{dat}&`tell sb.'&(\ref{chap:22}:\ref{ex:text16-20})\\
        \textit{rova} \textit{kahi}&from&`disappear from'&(\ref{chap:26}:\ref{ex:text19-34})(\ref{ex:sketch-47})\\
        \textit{tei} \textit{me}&with&`have sth.' (lit. `be there with sth.')&(\ref{chap:26}:\ref{ex:text19-4})(\ref{ex:sketch-81b})\\
        \textit{kuu} \textit{gunaha} \textit{vo}&\textsc{goal}&`fall down to'&(\ref{ex:sketch-17})\\
    \lspbottomrule
\end{tabular}
}
\end{table}

But there are also examples of VCs with incorporated prepositions that are headed by an Adjective or a Noun, see \tabref{tab:vcadjnincorp}.

\begin{table}
    \caption{VCs headed by an Adjective or Noun with an incorporated preposition.} \label{tab:vcadjnincorp}
\begin{tabularx}{\textwidth}{lXll}
    \lsptoprule
    \textsc{text}&\textsc{gloss of prep}&\textsc{translation}&\textsc{ref}\\
    \midrule
        \textit{mataa} \textit{ki}&\textsc{dat}&`be good for sb.'&(\ref{chap:16}:\ref{ex:text11-7})\\
        \textit{mataa} \textit{suku}&because.of&`be good because of sth.'&(\ref{chap:11}:\ref{ex:text7-72})\\
        \textit{taba} \textit{beera} \textit{ki}&\textsc{dat}&`be an important thing for sb.'&(\ref{ex:sketch-31})\\
    \lspbottomrule
\end{tabularx}
\end{table}

\subsection{The TM-Markers signifying temporal relations and modality}\label{sec:tmmarker}
Eight prehead TM-Markers signify how the situation described by the VC is located in time with respect to some other past, present or future situation given by the context and its modality. The TM-Markers are listed below with short meaning descriptions, and their frequencies in this text collection in square brackets, see \tabref{tab:tmmarkers}.

\begin{table}[b]
    \caption{The TM-markers and their meanings.} \label{tab:tmmarkers}
\begin{tabularx}{\textwidth}{llQl}
    \lsptoprule
    \textsc{tm}&\textsc{gloss}&\textsc{name and meaning}&\textsc{tokens}\\
    \midrule
%     \endfirsthead
%     \multicolumn{4}{c}%
%     {\tablename\ \thetable{} -- continued from previous page}\\
%     \lsptoprule
%     \textsc{tm}&\textsc{gloss}&\textsc{name and meaning}&\textsc{tokens}\\
% 	    \midrule
%     \endhead
%     \hline \multicolumn{4}{r}{{Continued on next page}}\\
%     \endfoot
%     \endlastfoot
    
        \textit{paa}&\textsc{tm}3&Change-of-Situation Marker: This situation followed, will follow, or generally follows another situation.&452 tokens\\
        \textit{na}&\textsc{tm}2&Realis Marker: This situation is a fact that occurred before the utterance or before some other past situation. The construction \textsc{tm}2-Head-\textsc{ipfv} means that the situation is or was simultaneous with the present or a past situation.&234 tokens\\
        \textit{toro}&\textsc{tm}7&Necessity Marker: This situation is/was necessary, obligatory or inevitable.&86 tokens\\
        \textit{kahi}&\textsc{tm}1&Posterior Marker1: This situation is/was expected to occur.&41 tokens\\
        \textit{pasi}&\textsc{tm}4&Posterior Marker2: This situation is or was wanted to occur.&33 tokens\\
        \textit{tau}&\textsc{tm}5&Posterior Marker3: This situation is or was about to happen.&23 tokens\\
        \textit{mene}&\textsc{tm}6&Timitive Marker:\footnote{The adjective ``timitive'' has been created by Palmer ``for fears'' (\citeyear[133-134]{Palmer:2001cj}) and is used by Lichtenberk in his article ``Modality and Mood in Oceanic languages,'' (\citeyear[346--349]{Lichtenberk:2016ph}).} This frightening situation might happen or might have happened.&9 tokens\\
        \textit{pari}&\textsc{tm}8&Counterfactual Marker: This situation almost happened.\footnote{In the \textsc{mted} the TM-Marker \textit{pari} is classified as an Adverb.}&no tokens\\
    \lspbottomrule
\end{tabularx}
\end{table}


The numbers of the glosses \textsc{tm}1--\textsc{tm}8 do not have any meaning as they were given to the particles at the beginning of the Teop language project.

The frequency and distribution of the TM-Markers in this text collection is obviously influenced by its genre, as most clauses make generic statements about plants and the stages of work processes. Even the making of the turtle fishing net that hasn't been practiced since the 1980s is described as a timeless practice. Statements about specific events of the past are rare.

The Realis Marker \textit{na} \textsc{tm}2 occurs 234 times, but only four times in VCs  without the \textsc{ipfv} Marker. Two of them are about past practices in (\ref{ex:sketch-1}) and (\ref{ex:sketch-36}). The third one is about a single past event, see (\ref{ex:sketch-37}).
\newpage
\ea\label{ex:sketch-37}
    O karaa vai o hiki toutou na deedee maa rohori Solomon\\
    \gll o	karaa	vai	o	hiki	toutou	na	dee{$\sim$}dee	maa	roho	ori	solomon\\
     \textsc{art}3.\textsc{sg}	pandanus	\textsc{dem}6	\textsc{art}3.\textsc{sg}	not.have	thorns	[\textsc{tm}2	\textsc{red}{$\sim$}carry	hither	before]\textsubscript{VC}	3\textsc{pl}.\textsc{pron}	Solomon\\
    \glt `The pandanus without thorns was brought from the Solomon Islands.' (\ref{chap:13}:\ref{ex:text9-17})
\z

\newpage
The fourth one occurs in a procedural text in a clause about an action that was done before the current referent time, see (\ref{ex:sketch-38}).

\ea\label{ex:sketch-38}
    ei kou na kaku pisi momohu ara.\\
    \gll {ei kou} na kaku pisi momohu ara\\
    because \textsc{tm}2 break fold first	1\textsc{pl}.\textsc{in}.\textsc{pron}\\
    \glt `(We put the bamboo strip on the ground, then we open a sago-palm leaf,) because we had folded it before.' (\ref{chap:17}:\ref{ex:text12-23})
\z

Past events are also expressed by non-verbal predicates, by VCs without any TM-Marker, and by VCs that are marked by any of the other TM-Markers as, for example, the Change of Situation Marker \textit{paa} \textsc{tm}3, which has the highest frequency of TM-Markers, see (\ref{ex:sketch-39}).

\ea\label{ex:sketch-39}
    \ea\label{ex:sketch-39a}
        Nabunuu roho enam ahiki to kara bao.\\
        \gll nabunuu roho e-nam ahiki to {kara bao}\footnotemark\\
        in.the.past before \textsc{h}-1\textsc{pl}.\textsc{ex}.\textsc{pron} not.exist \textsc{nspec}3.\textsc{sg} grass\\
        \glt `(Grass is the food of cows.) In the past we did not have any grass.'
    \ex\label{ex:sketch-39b}
        O burumakao he nomaa koa\\
        \gll o burumakao he nomaa koa\\
        \textsc{art}2.\textsc{pl} cattle \textsc{advers} come just\\
        \glt `But the cows came,'
    \ex\label{ex:sketch-39c}
        meori paa ma neba bata kiri bari bono kara bao.\\
        \gll me=ori paa ma neba bata ki=ri bari bono {kara bao}\\
        and4=3\textsc{pl}.\textsc{pron}	[\textsc{tm}3 come plant \textsc{cont} \textsc{dat}=3\textsc{pl}.\textsc{objm}]\textsubscript{VC}	4.\textsc{pl}.\textsc{pron} \textsc{obj}.\textsc{art}3.\textsc{sg} grass\\
        \glt `and they planted the grass for them.' (\ref{chap:11}:\ref{ex:text7-12}-\ref{ex:text7-13})
    \z
\z\footnotetext{For subordinate clauses functioning as predicates see \sectref{sec:depclaus}.}

A similar sequence of clauses about the past is Naphtali Maion's text about the past practice of \textit{kaku} \textit{goroto} `butchering turtles' in (\ref{chap:26}:\ref{ex:text19-31}-\ref{ex:text19-33}).

The following text excerpt illustrates the changes of temporal relations and modality in four subsequent clauses. The text describes what a man should do when he has carved the hull of a canoe in the jungle and now wants some men to pull it to the village, see (\ref{ex:sketch-40}).

\newpage
\ea\label{ex:sketch-40}
    \ea\label{ex:sketch-40a}
        Evehee ean toro paku a sii tabaan\\
        \gll evehee e-an toro paku a sii tabaan\\
        but	\textsc{h}-2\textsc{sg}.\textsc{pron} \textsc{tm}7 make \textsc{art}2.\textsc{sg} \textsc{dim} food\\
        \glt `(..., you tell some people to assist you to pull it home. This work does not cost anything,) But you should make some food'
    \ex\label{ex:sketch-40b}
        to kahi ma an koa vue o manae\\
        \gll to kahi ma an koa vue o manae\\
        \textsc{rel} \textsc{tm}1 come eat just particulary \textsc{art}2.\textsc{pl} young.men\\
        \glt `that the young men will eat,'
    \ex\label{ex:sketch-40c}
        to paa nao nau tea rahi sinivi.\\
        \gll to paa nao na=u tea rahi sinivi\\
        \textsc{rel} \textsc{tm}3 go thither=\textsc{pfv} \textsc{compl}1 pull canoe\\
        \glt `who have gone there to pull the canoe.' (\ref{chap:24}:\ref{ex:text18-36})
    \z
\z

Each of the four clauses of (\ref{ex:sketch-40}) has a distinct reference time that is given by the preceding context. In (\ref{ex:sketch-40c}) the Change-of-Situation Marker \textit{paa} (\textsc{tm}3) signifies that the situation of going (into the jungle) follows the situation of `you tell some people', whereas the \textsc{pfv}-Marker signifies that it precedes the situation of pulling the canoe.

The next example illustrates the use of the two Posteriority Markers \textit{tau} \textsc{tm}5 `be about to', \textit{pasi} \textsc{tm}4, and again the Necessity Marker \textit{toro} \textsc{tm}7 `must'. The Posteriority Marker2 \textit{pasi} (\textsc{tm}4) is combined with the \textsc{ipfv} Marker to signify habituality.\footnote{There is no example of \textit{pasi} \textsc{tm}4 without an aspect marker in this text collection. For the use of \textit{pasi} without the \textsc{ipfv} Marker and further information see \url{https://dictionaria.clld.org/units/teop-pasi}.} As in (\ref{ex:sketch-40}) the reference times (RefTs) change from clause to clause, see (\ref{ex:sketch-41}).

\ea\label{ex:sketch-41}
    \ea\label{ex:sketch-41a}
        Beara tau nao tea suvasuva,\\
        \gll be=ara tau nao tea suva{$\sim$}suva\\
        when=1\textsc{pl}.\textsc{in}.\textsc{pron} \textsc{tm}5 go \textsc{compl}1 \textsc{red}{$\sim$}poison\\
        \glt `Before we go to poison fish,'
    \ex\label{ex:sketch-41b}
        eara toro matatopo amaa taba\\
        \gll e-ara toro matatopo a=maa taba\\
        \textsc{h}-1\textsc{pl}.\textsc{in}.\textsc{pron} \textsc{tm}7 prepare \textsc{art}2.\textsc{sg}=\textsc{plm}	taba\\
        \glt `we must prepare the things'

\newpage
    \ex\label{ex:sketch-41c}
        to pasi dee raara.\\
        \gll to pasi dee ra=ara\\
        \textsc{rel} \textsc{tm}4 take 1\textsc{pl}.\textsc{in}.\textsc{ipfv}=1\textsc{pl}.\textsc{in}.\textsc{pron}\\
        \glt `that we'll take there.' (\ref{chap:37}:\ref{ex:text30-3})   
    \z
\z

The use of the Timitive Marker \textit{mene} \textsc{tm}6 is illustrated by a VC headed by \textit{hiki} `not have', see (\ref{ex:sketch-42}).

\ea\label{ex:sketch-42}
    \ea\label{ex:sketch-42a}
        Meara toro kikira nio naosinivi,\\
        \gll me=ara toro kikira ni=o nao-sinivi\\
        and4=1\textsc{pl}.\textsc{in}.\textsc{pron} \textsc{tm}7 look.after \textsc{appl}=\textsc{art}3.\textsc{sg} tree-canoe\\
        \glt `We must look after the canoe tree,'
    \ex\label{ex:sketch-42b}
        eara mene hiki sinivi.\\
        \gll e-ara mene hiki sinivi\\
        \textsc{h}-1\textsc{pl}.\textsc{in}.\textsc{pron} \textsc{tm}6 not.have canoe\\
        \glt `otherwise we would not have canoes.' (\ref{chap:7}:\ref{ex:text3-16})
    \z
\z

\subsection{The Aspect Markers}\label{sec:aspekt}

The Imperfective-Marker (\textsc{ipfv}) and the Perfective-Marker (\textsc{pfv}) hold the last position of the VC and exclude each other. The \textsc{ipfv} Marker inflects for person and number, see \tabref{tab:ipfvpfv}.


\begin{table}
    \caption{The Imperfective and the Perfective Aspect Markers and their meanings.} \label{tab:ipfvpfv}
\begin{tabularx}{\textwidth}{llX}
    \lsptoprule
    \textsc{aspect marker}&\textsc{gloss}&\textsc{meaning}\\
    \midrule
%     \endfirsthead
%     \multicolumn{3}{c}%
%     {\tablename\ \thetable{} -- continued from previous page}\\
%     \lsptoprule
%     \textsc{aspect marker}&\textsc{gloss}&\textsc{meaning}\\
% 	    \midrule
%     \endhead
%     \hline \multicolumn{3}{r}{{Continued on next page}}\\
%     \endfoot
%     \endlastfoot
        \textit{nom}&\textsc{ipfv}&\multirow{4}{7cm}{The situation is timeless, or it lasts for some time and includes the utterance time or the situation time of some other situation before or after the utterance time.}\\
        \textit{nana}, =\textit{na}, \textit{na}=&3\textsc{sg}.\textsc{ipfv}&\\
        \textit{rara}, =\textit{ra}, \textit{ra}=&1\textsc{pl}.\textsc{in}.\textsc{ipfv}&\\ 
        \textit{rori}, =\textit{ri}, \textit{ri}=&3\textsc{pl}.\textsc{ipfv}&\\\midrule
        =\textit{(v)u}, =\textit{v}=&\textsc{pfv}\footnote{In the \textsc{mted} the Perfective Aspect Marker (\textsc{pfv}) is called Immediateness Marker and glossed by \textsc{imm}, because we thought that this marker only signified that a situation that is relevant for an immediately subsequent situation.}&The situation is, was, will or could be completed and is relevant to a subsequent situation.\\
    \lspbottomrule
\end{tabularx}
\end{table}


The \textsc{ipfv} Marker indexes the subject NP or Pronoun as in (\ref{ex:sketch-41c}) unless a 1\textsuperscript{st} or 2\textsuperscript{nd} person Pronoun functions as the object and a NP or 3\textsuperscript{rd} person Pronoun as the subject. Then it indexes the object Pronoun, see (\ref{ex:sketch-43}).

\ea\label{ex:sketch-43}
    me eve kahi ani vakavakavara kurusu araraara\\
    \gll me eve kahi ani va-kava{$\sim$}kavara kurusu ara=ra=ara\\
    and4 3\textsc{sg}.\textsc{pron} \textsc{tm}1 eat \textsc{caus}-\textsc{red}{$\sim$}finished indeed 1\textsc{pl}.\textsc{in}.\textsc{objm}=1\textsc{pl}.\textsc{in}.\textsc{ipfv}=1\textsc{pl}.\textsc{in}.\textsc{pron}\\
    \glt `and he is going to eat us all.' (0513\_DoReCo\_doreco\_teop1238\_Sii\_06, in \textcite{Mosel:2024yp})
\z

The 3\textsuperscript{rd} person plural \textsc{ipfv} Marker regularly indexes subject plural Pronouns and NPs with human referents as in (\ref{ex:sketch-9}), (\ref{ex:sketch-14}), but sometimes also things and animals that play an important role in their context, see (\ref{chap:30}:\ref{ex:text23-8}).

The form \textit{nom} indexes 1\textsc{sg}, 1\textsc{pl}.\textsc{ex}., 2\textsc{sg}, and 2\textsc{pl} subject Pronouns, but is simply glossed by \textsc{ipfv}, because it may also index 1\textsc{pl}.\textsc{in}, 3\textsc{sg} and 3\textsc{pl} arguments, if the pre-head TM marker is not the Realis Marker \textit{na} \textsc{tm}2, see (\ref{ex:sketch-44}).

\ea\label{ex:sketch-44}
    A bua kehaa bona to kahi ato vamasuu bata ni nom bona kave\\
    \gll a	bua	kehaa	bona	to	kahi	ato	va-masuu	bata	ni	nom	bona	kave\\
    \textsc{art}2.\textsc{sg}	two2.\textsc{sg}\footnotemark{}	shell	\textsc{ana}.\textsc{sg}	\textsc{rel}	[\textsc{tm}1	hold	\textsc{advr}-steady	\textsc{cont}	\textsc{appl}	\textsc{ipfv}]\textsubscript{VC}	\textsc{obj}.\textsc{art}2.\textsc{sg}	net\\
    \glt `These two shells (i.e. heavy helmet shells) will hold the net steady.' (\ref{chap:27}:\ref{ex:text20-29})
\z\footnotetext{See: \url{https://dictionaria.clld.org/units/teop-bua}.}

The clitic forms of the \textsc{ipfv} and \textsc{pfv} Markers are attached to words preceding them within the VC or to a word directly following them beyond the VC.\footnote{Search (Words extra) in the \textsc{mted} for the Gloss ``IPFV'' and find 84 entries. Then search for the Gloss ``IMM'' and find 46 entries for the \textsc{pfv} Marker.}

TM Markers combine with Aspect Markers in various ways, but they may also occur without any Aspect Markers, and \textsc{ipfv} Markers may lack a TM marker, see \tabref{tab:distipfvpfv}.


\begin{table}
    \caption{The distribution of TM Markers Aspect Markers.} \label{tab:distipfvpfv}
\fittable{
\begin{tabular}{lllll}
    \lsptoprule
    &\textsc{gloss}&\textsc{v}&\textsc{v - ipfv}&\textsc{v - pfv}\\
    \midrule
        \textit{paa}&\textsc{tm}3&(\ref{ex:sketch-8}), (\ref{ex:sketch-12}), (\ref{ex:sketch-17}), (\ref{ex:sketch-27b}), (\ref{ex:sketch-39c})&(\ref{chap:17}:\ref{ex:text12-44}), (\ref{chap:24}:\ref{ex:text18-38})&(\ref{ex:sketch-40c}), (\ref{ex:sketch-47}), (\ref{chap:27}:\ref{ex:text20-20})\\
        \textit{na}&\textsc{tm}2&(\ref{ex:sketch-36}), (\ref{ex:sketch-37}), (\ref{ex:sketch-38})&(\ref{ex:sketch-1}), (\ref{ex:sketch-9}), (\ref{ex:sketch-11}), (\ref{ex:sketch-13})-(\ref{ex:sketch-15})&-\\
        \textit{toro}&\textsc{tm}7&(\ref{ex:sketch-40a}), (\ref{ex:sketch-41b}), (\ref{ex:sketch-42a})&(\ref{chap:5}:\ref{ex:text01-151})&-\footnote{In other texts \textit{tau} \textsc{tm}4, \textit{toro} \textsc{tm}7 and \textit{mene} \textsc{tm}6 are found in VCs, but they precede the \textsc{tm}3-V-\textsc{pfv}-Construction, e.g. \textit{enaa} \textit{tau} \textit{paa} \textit{gono=u} `I was about to get ...'}\\
        \textit{kahi}&\textsc{tm}1&(\ref{ex:sketch-40b})&(\ref{ex:sketch-43}), (\ref{ex:sketch-44})&(\ref{ex:sketch-48}), (\ref{chap:30}:\ref{ex:text23-15})\\
        \textit{pasi}&\textsc{tm}2&other texts\footnote{See \url{https://dictionaria.clld.org/units/teop-pasi}.}&(\ref{ex:sketch-41c}), (\ref{chap:17}:\ref{ex:text12-19})&(\ref{ex:sketch-49}), (\ref{chap:35}:\ref{ex:text28-8})\\
        \textit{tau}&\textsc{tm}5&(\ref{ex:sketch-41a})&(\ref{ex:sketch-35})&-\\
        \textit{mene}&\textsc{tm}6&(\ref{ex:sketch-42b})&(\ref{chap:27}:\ref{ex:text20-30})&-\\
        $\varnothing$&\textsc{tm}&(\ref{ex:sketch-39b}), (\ref{chap:19}:\ref{ex:text14-17})&(\ref{ex:sketch-45}), (\ref{ex:sketch-46}), (\ref{chap:6}:\ref{ex:text2-5})&-\\
    \lspbottomrule
\end{tabular}
}
\end{table}


The VC of relative clauses obligatorily lacks the Realis Marker \textit{na} \textsc{tm}2, so that only the \textsc{ipfv} Marker marks it to express habitual situations, see (\ref{ex:sketch-45}).


\newpage
\ea\label{ex:sketch-45}
    A rupi a bero kurusu to vahuhu nana a goroto.\\
    \gll a	rupi	a	bero	kurusu	to	vahuhu	nana	a	goroto\\
    \textsc{art}3.\textsc{pl}	egg	\textsc{art}3.\textsc{pl}	many	very	\textsc{rel}	lay.eggs	3\textsc{sg}.\textsc{ipfv}	\textsc{art}2.\textsc{sg}	turtle\\
    \glt `The turtle lays many eggs indeed.' (\ref{chap:26}:\ref{ex:text19-16})
\z

Further types of clause in which the VC obligatorily lacks the Realis Marker \textit{na} \textsc{tm}2, but may be marked for the imperfective aspect are the clauses with the Double Negator (\ref{chap:5}:\ref{ex:text01-6}), Adverbial Clauses (\ref{chap:6}:\ref{ex:text2-5}), Complement clauses (\ref{chap:5}:\ref{ex:text01-9}), and Imperative Clauses, see (\ref{ex:sketch-46}).


\ea\label{ex:sketch-46}
    Kikira vakikis bata koa ni nom eve.\\
    \gll kiikira	va-kikis	bata	koa	ni	nom	eve\\
    watch	\textsc{advr}-strong	\textsc{cont}	just	\textsc{appl}	\textsc{ipfv}	3\textsc{sg}.\textsc{pron}\\
    \glt `Just keep carefully watching it.' (\ref{chap:35}:\ref{ex:text28-18})
\z

The \textsc{pfv}-Marker combines most frequently with the Change-of-Situation Marker \textit{paa} \textsc{tm}3. If the reference time is simultaneous with  utterance time, the clause refers to a past event that is still relevant, see (\ref{ex:sketch-47}).

\ea\label{ex:sketch-47}
    Nomana amaa matapaku vai paa rova kahi avuaraara.\\
    \gll nomana	a=maa	matapaku	vai	paa	rova	kahi	a=vu=araara\\
    today	\textsc{art}2.\textsc{sg}=\textsc{plm}	practice	\textsc{dem}6	\textsc{tm}3	disappear	from	\textsc{objm}=\textsc{pfv}=1\textsc{pl}.\textsc{in}.\textsc{obj}.\textsc{pron}\\
    \glt `Today these practices have disappeared from us.' (\ref{chap:26}:\ref{ex:text19-34})
\z

When the \textsc{pfv}-Marker combines with the Posteriority Marker \textit{kahi} \textsc{tm}1, the event is located after the RefT and relevant for a following event, see (\ref{ex:sketch-48}).

\ea\label{ex:sketch-48}
    me kahi huha vuara bona maa bon amaa bero.\\
    \gll me	kahi	huha	vu=ara	bona	maa	bon	a=maa	bero\\
    and4	\textsc{tm}1	tie	\textsc{pfv}=1\textsc{pl}.\textsc{in}.\textsc{pron}	\textsc{obj}.\textsc{art}2.\textsc{sg}	\textsc{plm}	day	\textsc{art}2.\textsc{sg}=\textsc{plm}	many\\
    \glt `(The work is big,) and we will tie (the net) many days.' (\ref{chap:27}:\ref{ex:text20-23})
\z

The same kind of temporal relation is signified by the Posteriority Marker \textit{pasi} \textsc{tm}4 in combination with the \textsc{pfv}-Marker, see (\ref{ex:sketch-49}).

\ea\label{ex:sketch-49}
    O dorana o naono pasi vaasun vori,\\
    \gll o	dorana	o	naono	pasi	vaa-sun	v=ori\\
    \textsc{art}3.\textsc{sg}	four	\textsc{art}3.\textsc{sg}	pole	\textsc{tm}4	\textsc{caus}-stand	\textsc{pfv}=3\textsc{pl}.\textsc{pron}\\
    \glt `They will put up four poles, (and then they lay another four poles across, ...)' (\ref{chap:6}:\ref{ex:text2-34})
\z

\section{The Noun Phrase (NP)}\label{sec:np}
\subsection{Introduction}\label{sec:npintro}

NPs are marked by an Article which holds the first position of the phrase and is only dropped in the fixed expressions, e.g. \textit{meha} \textit{maa} \textit{vuri} `sometimes' (lit. `other times'), \textit{meha} \textit{papana} `other people, others' (lit. `other group') at the beginning of a clause.\footnote{See (\ref{chap:5}:\ref{ex:text01-83}), (\ref{chap:5}:\ref{ex:text01-144}).} NPs are headed by a Noun, a Verb, and rarely by an Adjective. The head may be preceded by additional determiners and followed by other content words. The last position of the NP is optionally held by the Anaphoric Demonstrative or another kind of Demonstrative. The sequential order of the constituents of basic NPs is shown in (\ref{ex:sketch-50}).

\ea\label{ex:sketch-50}
    \textup{[\textsc{art} determiners specifiers HEAD X \textsc{ana}/\textsc{dem}]\textsubscript{NP}}
\z

The Article signifies the lexical class of the head of the NP, which is different from its word class, and the specificity and grammatical number category of the NP. The paradigm of the Articles is presented in \sectref{sec:paraarticles}, and the classification of NP heads is described in \sectref{sec:classnphead}. Then \sectref{sec:preheadnp} deals with pre-head determiners and other words that follow the Article and further specify the reference of the NP.

The post-head position X can be held by a modifying content word, see \sectref{sec:modjux}, a modifying LnkP (\sectref{sec:subphrase}) or a content word functioning as the complement of a relational NP head, see \tabref{tab:relnoun}.


\begin{table}
    \caption{Relational Nouns and Verbs with complements in NPs.} \label{tab:relnoun}
\begin{tabularx}{\textwidth}{llll}
    \lsptoprule
    &\textsc{text}&\textsc{translation}&\textsc{ref}\\
    \midrule
        \textsc{art} N V&\textit{o kaku vagana}&`a kind of fishing'&(\ref{ex:sketch-10})\\
        \textsc{obj}.\textsc{art} N N&\textit{bene ta boon}&`piece of mangrove wood'&(\ref{ex:sketch-13})\\
        \textsc{art} V N&\textit{a deedee overe}&`the carrying of coconuts'&(\ref{chap:5}:\ref{ex:text01-157})\\
    \lspbottomrule
\end{tabularx}
\end{table}


NPs can be preceded by an attributive NumP and followed an attributive NP, AdjP, or PrepP, see \sectref{sec:subphrase}. If the NP head is marked by the Possessive Marker, it is followed by a possessor NP or Pronoun, see \sectref{sec:posscon}.

NPs preceded by NumP or containing the pre-head Numeral \textit{bua}/\textit{bua} `two' are grammatically singular, see \sectref{sec:preheadnp}. Furthermore, NPs and Pronouns referring to inanimate objects are mostly grammatically singular forms.

\subsection{The paradigm of Articles}\label{sec:paraarticles}

The paradigm of Teop Articles distinguishes the three classes of NP heads, singular and plural of the NP and its specific or non-specific reference, see \tabref{tab:paradigmart}.


\begin{table}
\small
    \caption{Relational Nouns and Verbs with complements in NPs.} \label{tab:paradigmart}
\begin{tabularx}{\textwidth}{llQQllll}
    \lsptoprule
    &\textsc{gloss}&\multicolumn{2}{c}{\textsc{art}1}&\multicolumn{2}{c}{\textsc{art}2}&\multicolumn{2}{c}{\textsc{art}3}\\
    \cmidrule(lr){3-4}\cmidrule(lr){5-6}\cmidrule(lr){7-8}
    &&\textsc{sg}&\textsc{pl}&\textsc{sg}&\textsc{pl}\footnote{For the ``inverse number marking'' of 2\textsuperscript{nd} class Nouns and 3\textsuperscript{rd} class Nouns by Articles see \textcite[159--166]{Corbett:2000gy} who also discusses this number marking in Teop (pp. 164-165).}&\textsc{sg}&\textsc{pl}\\
    \midrule
    Specific Basic Article&\textsc{art}&\textit{e}&\textit{ere}, \textit{o}&\textit{a}&\textit{o}&\textit{o}&\textit{a}\\
    Specific Object Article&\textsc{obj}.\textsc{art}&\textit{bene}\footnote{The
        Article \textit{bene} is also used with reflexive pronouns, search the \textsc{mted} Headword for ``{\^{}}bene''.
        }, \textit{bone}&\textit{bere}, \textit{benere}&\textit{bona}&\textit{bono}&\textit{bono}&\textit{bona}\\
    Non-specific Article&\textsc{nspec}&-&-&\textit{ta}&\textit{to}&\textit{to}&\textit{ta}\\
    Partitive Article&\textsc{part}.\textsc{art}& \textit{sa}\footnote{See: \url{https://dictionaria.clld.org/units/teop-sa_1}.} & sa& sa& sa& sa& sa\\
    \lspbottomrule
\end{tabularx}
\end{table}


If the Specific Article precedes a NP in a non-predicative function, the NP refers to an entity that is identifiable by the speech situation or the context. The Specific Object Article marks a non-topical object NP in clauses with a 3\textsuperscript{rd} person topical argument, see \sectref{sec:verbal}.

A predicative NP with a Specific Article states that its referent is identical with the referent of the subject NP or Pronoun, see (\ref{ex:sketch-51}),

\ea\label{ex:sketch-51}
    O haopi a hena na iana ...\\
    \gll o	haopi	a	hena	n=a	iana\\
    \textsc{art}3.\textsc{sg}	wrasse	\textsc{art}2.\textsc{sg}	name	3\textsc{sg}.\textsc{poss}=\textsc{art}2.\textsc{sg}	fish\\
    \glt `Haopi (Wrasse) is the name of a fish (that stays underground).' (\ref{chap:35}:\ref{ex:text28-1})
\z

\noindent or that the subject referent belongs to a specific kind of entity, see (\ref{ex:sketch-52}).

\ea\label{ex:sketch-52}
    A goroto a iana ...\\
    \gll a	goroto	a	iana\\
    \textsc{art}2.\textsc{sg}	turtle	\textsc{art}2.\textsc{sg}	swimming.animal\\
    \glt `The turtle is a swimming animal ...' (\ref{chap:26}:\ref{ex:text19-2})
\z

The plural form \textit{ere} \textsc{art}1.\textsc{pl} marks associative plurality, e.g. \textit{ere} \textit{iaa} `my mother and my aunts'\footnote{See: \url{https://dictionaria.clld.org/units/teop-ere}, also (\cite[66--67]{Mosel:2014gz}), (\cite[84, 101]{Corbett:2000gy}), and (\cite{Daniel:2013le}).}; it is not found in this text collection. The plural forms of the 2\textsuperscript{nd} and 3\textsuperscript{rd} class of specific and non-specific Articles denote groups of entities as, for example, the young breadfruit trees that are used for turtle nets, see (\ref{ex:sketch-53}).

\ea\label{ex:sketch-53}
    A naono baneo vai a rutarutaa na mataa nana, ...\\
    \gll a	naono	baneo	vai	a	ruta{$\sim$}rutaa	na	mataa	nana\\
    \textsc{art}3.\textsc{pl}	tree	breadfruit	\textsc{dem}6	\textsc{art}3.\textsc{pl}	\textsc{red}{$\sim$}small	\textsc{tm}2	good	3\textsc{sg}.\textsc{ipfv}\\
    \glt `The small breadfruit trees are good (because they are moist).' (\ref{chap:8}:\ref{ex:text4-12})
\z

Distributional plurality, i.e. the distribution of objects over space or time as in (\ref{ex:sketch-54}) and (\ref{ex:sketch-55}) below, is expressed by NPs headed by the Plural Marker \textit{maa}, see \sectref{sec:classnphead}.

The Non-Specific Article (\textsc{nspec}) marks NPs denoting entities whose existence or presence is denied, see (\ref{ex:sketch-54}),

\newpage
\ea\label{ex:sketch-54}
    Amaa hum bona kou na hiki nana to vasu, ...\\
    \gll a=maa	hum	bona	kou	na	hiki	nana	to	vasu\\
    \textsc{art}2.\textsc{sg}=\textsc{plm}	place	\textsc{ana}.\textsc{sg}	\textsc{emph}2	\textsc{tm}2	not.exist	3\textsc{sg}.\textsc{ipfv}	\textsc{nspec}2.\textsc{pl}	stone\\
    \glt `(And the turtle netting places ... are Taenaruene, Teopasuna, Teanin.) These places do not have any stones, (just white sand).' (\ref{chap:28}:\ref{ex:text21-15})
\z

\noindent or entities whose particular identity is not known yet, see (\ref{ex:sketch-55}).

\ea\label{ex:sketch-55}
    ta maa naono baneo vai ta maa rutarutaa.\\
    \gll ta	maa	naono	baneo	vai	ta	maa	ruta{$\sim$}rutaa\\
    \textsc{nspec}2.\textsc{sg}	\textsc{plm}	tree	breadfruit	\textsc{dem}6	\textsc{nspec}2.\textsc{sg}	\textsc{plm}	\textsc{red}{$\sim$}small\\
    \glt `(You must go and cut) some small breadfruit trees.' (\ref{chap:8}:\ref{ex:text4-7})
\z

The Partitive Article \textit{sa} `some' (\textsc{part}.\textsc{art}) signifies that the NP refers to some or any part of a whole, see (\ref{ex:sketch-56}).

\ea\label{ex:sketch-56}
    Na pakupaku raara bona kanono ae a iri ae sa taba sa ponisi mene dee gunaha bona kanono.\\
    \gll na	paku{$\sim$}paku	ra=ara	bona	kanono	ae	a	iri	ae	sa	taba	sa	ponisi	mene	dee	gunaha	bona	kanono\\
    \textsc{tm}2	\textsc{red}{$\sim$}do	1\textsc{pl}.\textsc{in}.\textsc{ipfv}=1\textsc{pl}.\textsc{in}.\textsc{pron}	\textsc{obj}.\textsc{art}2.\textsc{sg}	string	and1	\textsc{art}2.\textsc{sg}	hook	and1	\textsc{part}.\textsc{art}	thing	\textsc{part}.\textsc{art}	heavy	for	get	go.down	\textsc{obj}.\textsc{art}2.\textsc{sg}	string\\
    \glt `(We do the angling with a string and a hook) and something heavy to get the string down.' (\ref{chap:33}:\ref{ex:text26-2})
\z

The examples (\ref{ex:sketch-53}), (\ref{ex:sketch-55}) and (\ref{ex:sketch-56}) also illustrate that AdjPs are marked by an Article of the same kind of specificity, class and number as the NP they relate to as an attribute.

The Articles do not distinguish between definite and indefinite NPs. Indefiniteness can be optionally marked by the Indefiniteness Marker, see \sectref{sec:preheadnp}, and definiteness by the Anaphoric Demonstratives \textit{bona} (\textsc{ana}.\textsc{sg}) and \textit{bari} (\textsc{ana}.\textsc{pl}).

\subsection{The Articles and the classification of NP heads}\label{sec:classnphead}

Nouns, Verbs and Adjectives heading a NP are classified with respect to their compatibility with one or two of the three Article classes. Their classification shows some variation, which is determined by certain grammatical constructions or a speaker's individual choice. In the \textsc{mted} the classification of Nouns is indicated in the Part of Speech field.\footnote{For further information see \url{https://dictionaria.clld.org/contributions/teop\#section7} and (\cite{Mosel:2014kg}).} It is largely predictable:

\begin{itemize}
    \item N1: Proper names of people, kinship terms, Nouns denoting social roles, and domestic animals.\footnote{Due to the genre of this text collection, the Article of the 1\textsuperscript{st} class is only presented by a total of 14 tokens with singular NPs headed by \textit{beera} `chief', \textit{guu} `pig' and \textit{ta} `piece'.}
    \item N2: Nouns denoting kinds of human beings, vertebrates, invertebrates with legs, landmarks, food items, and parts of a whole; collective Nouns and four words denoting distinct kinds of plurality, see below.
    \item N3: Nouns denoting plants, parts of plants, and many things made of plant materials. Relational Nouns that are derelationalized by the suffix \textit{na} \textsc{derel}, see \sectref{sec:posscon}.
    \item N1/N2: animal names of the N2 class are shifted to the N1 class when the animals figure as protagonists in folk stories. The Noun \textit{ta} `piece', e.g. \textit{bene} \textit{ta} \textit{boon} `a piece of mangrove wood' (\ref{ex:sketch-13}) shifts to the N2 class, when it is used with a NumP or an AdjP, see \sectref{sec:nounta}.
    \item N2/N3: Nouns denoting a part of a whole, e.g. \textit{paka} `leaf', are used as an N2 in the Inalienable Possessive Construction and as a N3 in a simple \textsc{art}-N construction, see \sectref{sec:posscon}.
\end{itemize}

Verbs and Adjectives functioning as NP heads are, with two exceptions, affiliated with the N2 class (\cite[274--277, 279]{Mosel:2017as}). The exceptions are \textit{moroko} `speak' and \textit{upeupehe} `think'.

Four N2 words signify distinct kinds of plurality and form compounds with the following Noun referring to the quantified entity:

\largerpage
\begin{enumerate}
    \item \textit{maa} (\textsc{plm}), Plural Marker, signifies distributional plurality\footnote{Distributional plurality means that the NP refers to a plurality of entities that are distributed over space or time and do not form a group, cf. \textcite[11--117]{Corbett:2000gy}. Compare (\ref{ex:sketch-53}) with (\ref{ex:sketch-54}).}
    \item \textit{maamihu} `all sorts of', `every';
    \item \textit{ba} (\textsc{kin}.\textsc{plm}), the Kinship Plural Marker, denotes a group of people of the same generation sharing the same status in kin relations', e.g. kin relationships such as grandparents, mothers, etc.\footnote{See (\ref{ex:sketch-1}), (\ref{ex:sketch-37}), and \url{https://dictionaria.clld.org/units/teop-ba}.}
    \item \textit{tom} (\textsc{dyad}), the Dyadic Relationship Marker' means `group of two or more people who are in a reciprocal kin relationship', e.g. kin relationship such as a mother and her child.\footnote{Search (Words extra) in the \textsc{mted} for the Gloss ``DYAD'' and find 11 entries.} 
\end{enumerate}

The classification of these four words as N2-Nouns is based on the fact that they are always marked by the singular Article of the 2\textsuperscript{nd} class (\textsc{art}2.\textsc{sg}) irrespective of the inherent noun class of the following Noun which denotes the quantified entity. This is shown by NPs in which \textit{maa} \textsc{plm} is followed by an N3 Noun such as \textit{hum} `place' in (\ref{ex:sketch-54}) and \textit{naono} `plant', see (\ref{ex:sketch-57}).

\ea\label{ex:sketch-57}
    Amaa naono\\
    \gll a=maa	naono\\
    \textsc{art}2.\textsc{sg}=\textsc{plm}	plant\\
    \glt `the plants' (\ref{chap:10}:\ref{ex:text6-23})
\z

Similar to other Nouns, see \sectref{sec:modjux}, \textit{maa} \textsc{plm} forms compounds with Verbs, see (\ref{ex:sketch-58}):

\ea\label{ex:sketch-58}
    bona maa kiu vahihivaa te suburiori vai a otei.\\
    \gll bona maa kiu va-hihivaa te subu-riori vai a otei\\
    \textsc{obl}.\textsc{art}2.\textsc{sg} \textsc{plm} work \textsc{advr}-difficult \textsc{prep}.\textsc{art}1.\textsc{sg} grandparent-3\textsc{pl}.\textsc{poss}.\textsc{pron} \textsc{dem}6 \textsc{art}2.\textsc{sg} man\\
    \glt `the difficult work (lit. works) of their grandfather'\\\hfill (See \url{https://dictionaria.clld.org/units/teop-vahihivaa\_1})
\z

The N2 Noun \textit{maamihu} `all sorts of, every' is followed by N2-class Nouns, see (\ref{ex:sketch-59});

\ea\label{ex:sketch-59}
    a maamihu taba\\
    \gll a	maamihu	taba\\
    \textsc{art}2.\textsc{sg}	all.sorts.of	thing\\
    \glt `all the things' (\ref{chap:29}:\ref{ex:text22-4}), or `everything' (\ref{chap:33}:\ref{ex:text26-9})
\z

\noindent and by N3-class Nouns, see (\ref{ex:sketch-60}).
    
\ea\label{ex:sketch-60}
    a maamihu hum\\
    \gll a maamihu hum\\
    \textsc{art}2.\textsc{sg} all.sorts.of place\\
    \glt `all sorts of places, everywhere' (See \url{https://dictionaria.clld.org/units/teop-maamihu\_hum})
\z


\subsection{Pre-head words that specify the reference of the NP}\label{sec:preheadnp}
In addition to the Specific Article, the reference of the NP may be specified by the following six wordforms between the Article and the NP head.

\begin{enumerate}
    \item The Determiner called Indefiniteness Marker \textit{peha}\textsuperscript{2}, \textit{peho}\textsuperscript{2} (\textsc{indef}) `a, some, any' is a synonym of the Numeral \textit{peha}\textsuperscript{1}, \textit{peho}\textsuperscript{1} `one'; it is incompatible with the Numeral `one', but compatible with the Plural Marker \textit{maa}, see \sectref{sec:pluralmaa}.
    \item The Determiner \textit{meha}, \textit{meho} `some item or items (of a group), other' is incompatible with the Indefiniteness Marker and holds the same position.
    \item The Numeral \textit{peha}\textsuperscript{1}, \textit{peho}\textsuperscript{1} `one' is incompatible with the Indefiniteness Marker.
    \item The Numeral \textit{bua}, \textit{buo} `two' is compatible with the Indefiniteness Marker.
    \item The Diminutive Particle \textit{sii} (\textsc{dim}) means `little' or denotes the speaker's empathy.
    \item The Adjective \textit{vahara} `small (persons or things)' (small.\textsc{pl}).
\end{enumerate}

The sequential order of the pre-head elements is fixed, see \tabref{tab:prehead}.


\begin{table}
    \caption{The pre-head determiners in their sequential order in NPs.} \label{tab:prehead}
\begin{tabularx}{\textwidth}{llllllll}
    \lsptoprule
    \textsc{art}&\textsc{det}&\textsc{num}&\textsc{dim}&\textsc{adj}&\textsc{noun}&\textsc{translation}&\textsc{ref}\\
    \midrule
    \textit{A}&\textit{peha}&-&-&-&\textit{aba}&`a person, someone'&(\ref{chap:26}:\ref{ex:text19-18})\\
    \textit{O}&&\textit{peho}&&-&\textit{sinivi}&`one canoe'&(\ref{chap:30}:\ref{ex:text23-5})\\
    \textit{O}&\textit{meho}&\textit{buo}&&-&\textit{sinivi}&`the other two canoes'&(\ref{chap:30}:\ref{ex:text23-6})\\
    \textit{O}&&&\textit{sii}&\textit{vahara}&\textit{vasu}&`the little stones'&(\ref{chap:35}:\ref{ex:text28-6})\\
    \lspbottomrule
\end{tabularx}
\end{table}


The inflecting forms of the Numerals \textit{peha}\textsuperscript{1}, \textit{peho}\textsuperscript{1} `one', \textit{bua}, \textit{buo} `two' and the Determiners \textit{peha}\textsuperscript{2}, \textit{peho}\textsuperscript{2} \textsc{indef}, \textit{meha}, \textit{meho} `other' indicate the noun class and grammatical number of the NP. The Numeral \textit{peha}\textsuperscript{1}, \textit{peho}\textsuperscript{1} and the Determiner \textit{meha}/\textit{meho} may function as the head of anaphoric SubstPs, as illustrated by (\ref{ex:sketch-2}) and (\ref{ex:sketch-6}), see \sectref{sec:phrasefuncclause}.

\subsection{The Noun \textit{ta} `piece' and the \label{sec:nounta} Plural Marker \textit{maa}}\label{sec:pluralmaa}

The Nouns \textit{ta} and \textit{maa} have two properties in common:

\begin{enumerate}
    \item They head compounds and never occur by themselves.
    \item They are repeated in attributive AdjPs.
\end{enumerate}

The Noun \textit{ta} `piece' belongs to the N1/N2 class\footnote{See: \url{https://dictionaria.clld.org/units/teop-ta\_2}.} and is repeated in a modifying AdjP, see (\ref{ex:sketch-61}).

\ea\label{ex:sketch-61}
    A ta naono vai a ta beera, bebeahu ae kikis\\
    \gll a	ta	naono	vai	a	ta	beera	bebeahu	ae	kikis\\
    \textsc{art}2.\textsc{sg}	piece	wood	\textsc{dem}6	\textsc{art}2.\textsc{sg}	piece	big	long	and1	strong\\
    \glt `A big, long and strong pole' (\ref{chap:16}:\ref{ex:text11-3})
\z

Another Noun that is repreated in AdjP is \textit{kara} `vine', see (\ref{chap:11}:\ref{ex:text7-6}), (\ref{chap:11}:\ref{ex:text7-78}). The same repetition occurs optionally in NPs headed by \textit{maa} \textsc{plm}, see (\ref{ex:sketch-62}).

\ea\label{ex:sketch-62}
    tea maa nahu vai amaa beera\\
    \gll te=a	maa	nahu	vai	a=maa	beera\\
    \textsc{prep}=\textsc{art}2.\textsc{sg}	\textsc{plm}	pot	\textsc{dem}6	\textsc{art}2.\textsc{sg}=\textsc{plm}	large\\
    \glt `in large pots' (\ref{chap:6}:\ref{ex:text2-50})
\z

The Plural Marker \textit{maa} and \textit{ta} `piece' may form a complex compound with a third Noun\footnote{Some authors write \textit{mata} instead of \textit{maa} \textit{ta}, e.g. \textit{bona} \textit{mata} \textit{goroto} `the pieces of the turtle' (\ref{chap:26}:\ref{ex:text19-47}).}, see (\ref{ex:sketch-63}).

\ea\label{ex:sketch-63}
    Amaa ta naono bari mee amaa ta kakato koa.\\
    \gll a=maa	ta	naono	bari	mee	a=maa	ta	kakato	koa\\
    \textsc{art}2.\textsc{sg}=\textsc{plm}	piece	wood	\textsc{ana}.\textsc{pl}	also	\textsc{art}2.\textsc{sg}=\textsc{plm}	piece	short	quite\\
    \glt `These poles are quite short.' (\ref{chap:16}:\ref{ex:text11-24})
\z


\largerpage[2]
The repetition of the Plural Marker is also found with the Indefiniteness Marker, see (\ref{ex:sketch-64})

\ea\label{ex:sketch-64}
    Amaa peha maa vuri ...\\
    \gll a=maa	peha	maa	vuri\\
    \textsc{art}2.\textsc{sg}=\textsc{plm}	\textsc{indef}2.\textsc{sg}	\textsc{plm}	time\\
    \glt `Sometimes (they go diving for turtles).' (\ref{chap:26}:\ref{ex:text19-27})
\z

\noindent and the Determiner \textit{meha} `other', see (\ref{ex:sketch-65}).\clearpage

\ea\label{ex:sketch-65}
    Amaa meha maa ana\\
    \gll a=maa	meha	maa	ana\\
    \textsc{art}2.\textsc{sg}=\textsc{plm}	other2.\textsc{sg}	\textsc{plm}	root\\
    \glt `Other roots (have also started growing out (of the nut))' (\ref{chap:5}:\ref{ex:text01-17})
\z

\section{Pronouns}\label{sec:pronouns}
\subsection{The four person categories of Teop Pronouns}\label{sec:fourcat}

The Teop Pronouns distinguish between four persons, singular and plural, and five classes of morpho-syntactic forms: basic Pronouns (\textsc{pron}), Heavy Pronouns (\textsc{h}-\textsc{pron}), Object Pronouns (\textsc{obj}.\textsc{pron}), suffixed Possessive Pronouns (\textsc{poss}-\textsc{pron}), and the clitic Consecutive Pronoun (\textsc{pron}=), see \tabref{tab:pronouns}.


\begin{table}[b]
    \caption{The Teop Pronouns.} \label{tab:pronouns}
\begin{tabularx}{\textwidth}{XXXXXl}
    \lsptoprule
    &\textsc{pron}&\textsc{h-pron}&\textsc{obj.pron}&\textsc{poss-pron}&\textsc{pron=}\\
    \midrule
%     \endfirsthead
%     \multicolumn{6}{c}%
%     {\tablename\ \thetable{} -- continued from previous page}\\
%     \lsptoprule
%     &\textsc{pron}&\textsc{h-pron}&\textsc{obj.pron}&\textsc{poss-pron}&\textsc{pron=}\\
% 	    \midrule
%     \endhead
%     \hline \multicolumn{6}{r}{{Continued on next page}}\\
%     \endfoot
%     \endlastfoot
    1\textsc{sg}&\textit{naa}&\textit{e-naa}&\textit{a-naa}&$\varnothing$-\textit{naa}&\textit{ne}\footnote{See: \url{https://dictionaria.clld.org/units/teop-ne\_2}. The form occurs in the TLC by itself or follows a subject 1\textsc{sg}.\textsc{pron} or 2\textsc{sg}.\textsc{pron}, e.g. \textit{ne} \textit{nao} `then I go'. It does not occur in this text-collection.}\\
    1\textsc{pl}.\textsc{ex}&\textit{nam}&\textit{e-nam}&\textit{a-nam}&\textit{-ma-nam}&-\\
    1\textsc{pl}.\textsc{in}&\textit{ara}&\textit{e-ara}&\textit{ara-ara}&\textit{-ra-ara}&\textit{a=re}\\
    2\textsc{sg}&\textit{an}&\textit{e-an}&\textit{vu-an}&\textit{-m-an}&-\\
    2\textsc{pl}&\textit{am}&\textit{e-am}&\textit{ame-am}&\textit{-me-am}&-\\
    3\textsc{sg}&\textit{eve, e}&\textit{eove, eve}&\textit{eve, e}&\textit{-na-e}&\textit{o=re}\\
    3\textsc{pl}&\textit{ori}&\textit{e-ori}&\textit{ori}&\textit{-ri-ori}&\textit{o=re}\\
    4\textsc{sg}&-&-&\textit{bona}&-&-\\
    4\textsc{pl}&-&-&\textit{bari}&-&-\\
    \lspbottomrule
\end{tabularx}
\end{table}


The 1\textsuperscript{st} and 2\textsuperscript{nd} Person Pronouns differ morphologically and syntactically from the non-speech-act Pronouns which are called 3\textsuperscript{rd} and 4\textsuperscript{th} Person Pronouns. While the 1\textsuperscript{st} and 2\textsuperscript{nd} Person Object Pronouns are derived from the Basic Pronouns by prefixes, the 3\textsuperscript{rd} Person Pronouns do not distinguish between Basic and Object Pronouns, e.g. \textit{eve}, e 3\textsc{sg}.\textsc{pron}, \textit{ori} 3\textsc{pl}.\textsc{pron}. But corresponding to the differential object marking of NPs, see \sectref{sec:verbal}, the 3\textsuperscript{rd} Person Basic Pronouns only function as an object if the subject Pronoun refers to a speech-act participant, see (\ref{ex:sketch-66}).

\ea\label{ex:sketch-66}
    Beara tara ahe eve, ...\\
    \gll be=ara	tara	ahe	eve\\
    when=1\textsc{pl}.\textsc{in}.\textsc{pron}	see	recognize	3\textsc{sg}.\textsc{pron}\\
    \glt `When we spot it, (we must push the fishing-stick into it.)' (\ref{chap:35}:\ref{ex:text28-7})
\z

If a NP or a 3\textsuperscript{rd} Person Pronoun functions as the subject, the object is formed by the 4\textsuperscript{th} Person Pronoun, see (\ref{ex:sketch-67}).

\ea\label{ex:sketch-67}
    ... a peha aba na antee nana tea paku bona.\\
    \gll a	peha	aba	na	antee	nana	tea	paku	bona\\
    \textsc{art}2.\textsc{sg}	one2.\textsc{sg}	person	\textsc{tm}2	can	3\textsc{sg}.\textsc{ipfv}	\textsc{compl}1	do	4\textsc{sg}.\textsc{pron}\\
    \glt `... a single person can do it.' (\ref{chap:24}:\ref{ex:text18-6})
\z

The plural forms of 3\textsuperscript{rd} and 4\textsuperscript{th} Person Pronouns are used in strikingly different ways. The form \textit{ori} `they, them' (3\textsc{pl}.\textsc{pron}) is obligatory when the Pronoun refers to more than one person, see (\ref{ex:sketch-68}).

\ea\label{ex:sketch-68}
    ean repaa vaan riori, \\
    \gll e-an	re=paa	va-an	ri=ori\\
    \textsc{h}=2\textsc{sg}.\textsc{pron}	\textsc{consec}=\textsc{tm}3	\textsc{caus}-eat	3\textsc{pl}.\textsc{objm}=3\textsc{pl}.\textsc{pron}\\
    \glt `(When they have pulled the canoe home,) you will give them to eat, ...' (\ref{chap:24}:\ref{ex:text18-37})
\z

With non-human referents the 3\textsuperscript{rd} Person Singular form is selected, see (\ref{ex:sketch-69}).

\ea\label{ex:sketch-69}
    Bean kosi vakavara a paka taovu, ean paa kaakaru e bono pesupaava, ...\\
    \gll be=an	kosi	va-kavara	a	paka	taovu	e-an	paa	kaakaru	e	bono	pesupaava\\
     when=2\textsc{sg}.\textsc{pron}	cut.off.leaves	\textsc{advr}-complete	\textsc{art}3.\textsc{pl}	leaf	giant.pandanus	\textsc{h}-2\textsc{sg}.\textsc{pron}	\textsc{tm}3	scrape.thorns.off	3\textsc{sg}.\textsc{pron}	\textsc{obj}.\textsc{art}3.\textsc{sg}	coconut.husk\\
    \glt `After you have cut off the pandanus leaves, you scrape off the thorns with a coconut husk ...' (\ref{chap:14}:\ref{ex:text10-9})
\z

Exceptions are found in contexts where animals function as discourse topics as in the descriptions of marine animals. In other cases as in (\ref{ex:sketch-21}), the writer's usage of the plural form seems to be influenced by English. Consequently, \textit{e}, \textit{eve} \textsc{3sg}.\textsc{pron} is variably translated by English singular or plural pronouns of the 3\textsuperscript{rd} person.

The selection of \textit{bona} 4\textsc{sg}.\textsc{pron} or \textit{bari} 4\textsc{pl}.\textsc{pron} does not only depend on the number of their referents, but also on the number of the referent of the subject NP or Pronoun. The Pronoun \textit{bona} 4\textsc{sg}.\textsc{pron} is only used, when both the subject NP or Pronoun and \textit{bona}  4\textsc{sg}.\textsc{pron} refer to a single entity as in (\ref{ex:sketch-67}) `a single person can do it', or \textit{bona}  4\textsc{sg}.\textsc{pron} refers to inanimate entities.

The plural form \textit{bari} 4\textsc{pl}.\textsc{pron} refers, as expected, to a plurality of undergoers in clauses with a singular subject NP or Pronoun, see in (\ref{ex:sketch-70}).

\ea\label{ex:sketch-70}
    eove repaa mee ri bari rasuu tea aito.\\
    \gll eove	re=paa	mee	ri	bari	rasuu	tea	aito\\
    \textsc{h}.3\textsc{sg}.\textsc{pron}	\textsc{consec}=\textsc{tm}3	take	3\textsc{pl}.\textsc{objm}	4\textsc{pl}.\textsc{pron}	jungle	\textsc{compl}1	collect.herbs\\
    \glt `(A man who is an expert on herbs selects two or four men,) then he takes them into the jungle to collect herbs.' (\ref{chap:28}:\ref{ex:text21-3})
\z

But unexpectedly, \textit{bari} 4\textsc{pl}.\textsc{pron} also refers to a single undergoer in clauses with a subject NP or Pronoun referring to a plurality of actors, see (\ref{ex:sketch-71}).

\ea\label{ex:sketch-71}
    eori repaa tagava bari\\
    \gll e-ori	re=paa	tagava	bari\\
    \textsc{h}-3\textsc{pl}.\textsc{pron}	\textsc{consec}=\textsc{tm}3	fell	4\textsc{pl}.\textsc{pron}\\
    \glt `and then they fell it' (\ref{chap:6}:\ref{ex:text2-45})
\z

The form \textit{bari} 4\textsc{pl}.\textsc{pron} is not only used in clauses about present habitual events, but also in clauses about individual past events with more than one actor and a single undergoer, see (\ref{ex:sketch-72}).

\ea\label{ex:sketch-72}
    meori pae ni bari.\\
    \gll me=ori pae ni bari\\
    and4=3\textsc{pl}.\textsc{pron} bury \textsc{appl} 4\textsc{pl}.\textsc{pron}\\
    \glt `and they buried him.' (See \url{https://dictionaria.clld.org/units/teop-bari\_1})
\z

The reference to a plurality of undergoers in (\ref{ex:sketch-70}) and a single undergoer in (\ref{ex:sketch-71}) and (\ref{ex:sketch-72}) is not distinguished by the object Pronoun, but by the Object Marker \textit{ri} in the VC of (\ref{ex:sketch-70}).

The next four tables present an overview of the syntactic functions of the 1\textsuperscript{st} Person Inclusive Plural Pronouns, the 3\textsuperscript{rd}, the 4\textsuperscript{th} Person Pronouns, and zero anaphora. For reasons of space, the 1\textsuperscript{st} Person Inclusive Plural Pronoun was chosen as presentative for 1\textsuperscript{st} and 2\textsuperscript{nd} Person Pronouns, see \tabref{tab:firstincl}.


\begin{table}
    \caption{The functions of 1\textsuperscript{st} Person Inclusive Pronouns.} \label{tab:firstincl}
\begin{tabularx}{\textwidth}{lQQQQQ}
    \lsptoprule
    &\multicolumn{5}{c}{\textsc{syntactic functions}}\\
    \cmidrule{2-6}
    \textsc{form}&\textsc{clefted}&\textsc{topical}&\textsc{topical}&\textsc{non-top.}&\textsc{non-top.}\\
    \textsc{gloss}&\textsc{pron.}&\textsc{subj.}&\textsc{obj.}&\textsc{subj.}&\textsc{obj.}\\
    \midrule
        \textit{e-ara}&\multirow{2}{*}{+\footnote{There is no example in this text collection, but see \textit{Enaa} \textit{to} \textit{vigu} \textit{bona} \textit{sikuuru} \textit{bona}. `It was me who build this school.' (\cite[59]{Rum:2014jx}) and \url{https://dictionaria.clld.org/units/teop-vigu}.}}&\multirow{2}{*}{(\ref{chap:5}:\ref{ex:text01-53})}&\multirow{2}{*}{-}&\multirow{2}{*}{-}&\multirow{2}{*}{(\ref{ex:sketch-81a})}\\
        \textsc{h}-1\textsc{pl}.\textsc{in}.\textsc{pron}&&&&&\\\hline
        \textit{ara}&\multirow{2}{*}{-}&\multirow{2}{*}{-}&\multirow{2}{*}{-}&\multirow{2}{*}{(\ref{ex:sketch-10})}&\multirow{2}{*}{-}\\
        1\textsc{pl}.\textsc{in}.\textsc{pron}&&&&&\\\hline
        \textit{a}=&\multirow{2}{*}{-}&\multirow{2}{*}{(\ref{ex:sketch-78})}&\multirow{2}{*}{-}&\multirow{2}{*}{-}&\multirow{2}{*}{-}\\
        1\textsc{pl}.\textsc{in}.\textsc{pron}=&&&&&\\\hline
        \textit{araara}&\multirow{2}{*}{-}&\multirow{2}{*}{-}&\multirow{2}{*}{-}&\multirow{2}{*}{-}&\multirow{2}{*}{(\ref{chap:22}:\ref{ex:text16-20})}\\
        1\textsc{pl}.\textsc{in}.\textsc{obj}.\textsc{pron}&&&&&\\
    \lspbottomrule
\end{tabularx}
\end{table}


The syntactic functions of the 3\textsuperscript{rd} Person Pronouns differ from those of the 1\textsuperscript{st} and 2\textsuperscript{nd} Person Pronouns, see \tabref{tab:thirdpronfunc}.


\begin{table}
    \caption{3\textsuperscript{rd} Person Pronouns and their functions.} \label{tab:thirdpronfunc}
\begin{tabularx}{\textwidth}{lQQQQQ}
    \lsptoprule
    &\multicolumn{5}{c}{\textsc{syntactic functions}}\\
    \cmidrule{2-6}
    \textsc{form}&\textsc{clefted}&\textsc{topical}&\textsc{topical}&\textsc{non-top.}&\textsc{non-top.}\\
    \textsc{gloss}&\textsc{pron.}&\textsc{subj.}&\textsc{obj.}&\textsc{subj.}&\textsc{obj.}\\
    \midrule
        \textit{eove}&\multirow{2}{*}{(\ref{chap:20}:\ref{ex:text15-2})}&\multirow{2}{*}{(\ref{chap:5}:\ref{ex:text01-11})}&\multirow{2}{*}{-}&\multirow{2}{*}{-}&\multirow{2}{*}{(refl. obj.)\footnote{When functioning as a reflexive object, it is preceded by the Object Article \textit{bene} (\textsc{obj}.\textsc{art}1.\textsc{sg}), see  \url{https://dictionaria.clld.org/units/teop-eove}.}}\\
        \textsc{h}.3\textsc{sg}.\textsc{pron}&&&&&\\\hline
        \textit{eve}&\multirow{2}{*}{(\ref{chap:11}:\ref{ex:text7-7})}&\multirow{2}{*}{(\ref{chap:7}:\ref{ex:text3-8})}&\multirow{2}{*}{(\ref{chap:19}:\ref{ex:text14-2})}&\multirow{2}{*}{(\ref{chap:5}:\ref{ex:text01-146})}&\multirow{2}{*}{(\ref{chap:37}:\ref{ex:text30-10})\footnote{See: \url{https://dictionaria.clld.org/units/teop-eve}.}}\\
        3\textsc{sg}.\textsc{pron}&&&&&\\\hline
        =\textit{e}&\multirow{2}{*}{-}&\multirow{2}{*}{-}&\multirow{2}{*}{-}&\multirow{2}{*}{(\ref{chap:5}:\ref{ex:text01-165})}&\multirow{2}{*}{(\ref{chap:5}:\ref{ex:text01-141})}\\
        3\textsc{sg}.\textsc{pron}&&&&&\\\hline
        \textit{o}=&\multirow{2}{*}{-}&\multirow{2}{*}{(\ref{chap:5}:\ref{ex:text01-3})}&\multirow{2}{*}{-}&\multirow{2}{*}{-}&\multirow{2}{*}{-}\\
        3\textsc{sg}/\textsc{pl}.\textsc{pron}=&&&&&\\\hline
        \textit{e}-\textit{ori}&\multirow{2}{*}{+\footnote{\textit{Eori to vigu bona vaan vaa Kekesu}. `It was them who built the village of Kekesu.' (\cite[46]{Morekevan:2014es}).}}&\multirow{2}{*}{(\ref{chap:7}:\ref{ex:text3-11})}&\multirow{2}{*}{-}&\multirow{2}{*}{-}&\multirow{2}{*}{-}\\
        \textsc{h}-3\textsc{pl}.\textsc{pron}&&&&&\\\hline
        \textit{ori}&\multirow{2}{*}{-}&\multirow{2}{*}{(\ref{chap:7}:\ref{ex:text3-10})}&\multirow{2}{*}{-}&\multirow{2}{*}{(\ref{chap:7}:\ref{ex:text3-10})}&\multirow{2}{*}{(\ref{chap:24}:\ref{ex:text18-37})}\\
        3\textsc{pl}.\textsc{pron}&&&&&\\
    \lspbottomrule
\end{tabularx}
\end{table}


As illustrated by the examples (\ref{ex:sketch-67}), (\ref{ex:sketch-70}), (\ref{ex:sketch-71}) and (\ref{ex:sketch-72}), the distinction of singular and plural 4\textsuperscript{th} Person Pronouns is not exclusively determined by its antecedent, but also by the referent of the subject NP, so that depending on the context both the singular form \textit{bona} and the plural form \textit{bari} are translated by English singular or plural forms, see \tabref{tab:fourthperson}.


\begin{table}
    \caption{The 4\textsuperscript{th} person Pronouns and their English translations.} \label{tab:fourthperson}
\begin{tabularx}{\textwidth}{XXXl}
    \lsptoprule
    \textsc{pronoun}&\textsc{gloss}&\textsc{translation}&\textsc{ref}\\
    \midrule
        \textit{bona}&4\textsc{sg}.\textsc{pron}&it&(\ref{chap:5}:\ref{ex:text01-30}), (\ref{chap:17}:\ref{ex:text12-25}), (\ref{ex:sketch-67})\\
        &&them&(\ref{chap:10}:\ref{ex:text6-12}), (\ref{chap:28}:\ref{ex:text21-4})\\
        \textit{bari}&4\textsc{pl}.\textsc{pron}&it, him&(\ref{chap:5}:\ref{ex:text01-116}), (\ref{chap:7}:\ref{ex:text3-8}), (\ref{chap:26}:\ref{ex:text19-31}), (\ref{ex:sketch-71}), (\ref{ex:sketch-72})\\
        &&them&(\ref{chap:5}:\ref{ex:text01-91}), (\ref{chap:28}:\ref{ex:text21-3}, (\ref{ex:sketch-70})\\
    \lspbottomrule
\end{tabularx}
\end{table}


If an actor or undergoer of a situation is understood from the context, it may not be explicitly mentioned by a NP or Pronoun. This zero-anaphora is found in the topical position of subjects or direct/primary objects and in the position of non-topical direct/primary objects see \tabref{tab:zero}.


\begin{table}
    \caption{Zero anaphora.} \label{tab:zero}
\begin{tabularx}{\textwidth}{XXl}
    \lsptoprule
    &\textsc{subject}&\textsc{direct/primary object}\\
    \midrule
        Topical position&(\ref{chap:5}:\ref{ex:text01-7}), (\ref{chap:10}:\ref{ex:text6-11})&(\ref{chap:5}:\ref{ex:text01-43}), (\ref{chap:13}:\ref{ex:text9-15}), (\ref{chap:33}:\ref{ex:text26-2})\\
        Non-topical position&-&(\ref{chap:5}:\ref{ex:text01-154}), (\ref{chap:5}:\ref{ex:text01-155}), (\ref{chap:38}:\ref{ex:text31-4})\\
    \lspbottomrule
\end{tabularx}
\end{table}


The following sections describe the functions of the Heavy Pronoun in \sectref{sec:heavypron}, and the Consecutive Pronoun in \sectref{sec:consecpron}. For the Possessive Pronouns see \sectref{sec:posscon}. 

\newpage
\subsection{The Heavy Pronoun}\label{sec:heavypron}

The Heavy Pronoun functions as a clause initial topical subject, see (\ref{ex:sketch-73}),

\ea\label{ex:sketch-73}
    eove repaa taneo tea pura hopo vo kasuana.\\
    \gll eove	re=paa	taneo	tea	pura	hopo	vo	kasuana\\
    \textsc{h}.3\textsc{sg}.\textsc{pron}	\textsc{consec}=\textsc{tm}3	start	\textsc{compl}1	grow	enter	\textsc{goal}	ground\\
    \glt `(When the sprouting root comes out,) it starts growing into the ground,' (\ref{chap:5}:\ref{ex:text01-11})
\z

\noindent as a clause external focused Pronoun in the so-called ``it-cleft construction''\footnote{\textcite[133, 149, 152, 155]{Andrews:2007eq}; \textcite[17, 70--71]{Lambrecht:1994od}.}, see (\ref{ex:sketch-74}),

\ea\label{ex:sketch-74}
    Eove to kom, hio ae goroho raara.\\
    \gll eove	to	kom	hio	ae	goroho	ra=ara\\
    \textsc{h}.3\textsc{sg}.\textsc{pron}	\textsc{rel}	step.on	sit	and1	sleep	1\textsc{pl}.\textsc{in}.\textsc{ipfv}=1\textsc{pl}.\textsc{in}.\textsc{pron}\\
    \glt `(The floor is a thing in our house, the Teop people.) It is where we step, sit, and sleep on.' (\ref{chap:20}:\ref{ex:text15-2})
\z

\noindent as the head of a negated VC\footnote{Compare (\ref{ex:sketch-18}) in \sectref{sec:subphrase}.}, see (\ref{ex:sketch-75}),

\ea\label{ex:sketch-75}
    Saka enaa haana to ato.\\
    \gll saka e-naa haa=na to ato\\
    \textsc{neg}1 \textsc{h}-1\textsc{sg}.\textsc{pron} \textsc{neg}2=3\textsc{sg}.\textsc{ipfv} \textsc{rel} hold\\
    \glt `(It is) not me who took possession (of it).' (\cite[22]{Siarive:2014gg})
\z

\noindent and as a reflexive object Pronoun of a transitive Verb, see (\ref{ex:sketch-76}).

\ea\label{ex:sketch-76}
    menaa paa hivi komana enaa,\\
    \gll me=naa paa hivi komana e-naa\\
    and4=1\textsc{sg}.\textsc{pron} [\textsc{tm}3 ask self]\textsubscript{VC} \textsc{h}-1\textsc{sg}.\textsc{pron}\\
    \glt `and I asked myself,' (See \url{https://dictionaria.clld.org/units/teop-enaa})
\z

The 3\textsuperscript{rd} person reflexive Pronouns are preceded by the Article \textit{bene} \textsc{obj}.\textsc{art}1.\textsc{sg}, see (\ref{ex:sketch-77}).

\ea\label{ex:sketch-77}
    mepaa hivi komana bene eve, ...\\
    \gll me=paa hivi komana bene eve\\
    and4=\textsc{tm}3 red ask self \textsc{obj}.\textsc{art}1.\textsc{sg} 3\textsc{sg}.\textsc{pron}\\
    \glt `(a young boy found the vine of a yam) and he asked himself, ...' (See \url{https://dictionaria.clld.org/units/teop-bene\_eve})
\z

\subsection{The Consecutive Pronoun}\label{sec:consecpron}

The Consecutive Clause is a dependent non-embedded clause that is marked by the conjunction \textit{re} `then, so that' (\textsc{consec}) which holds the position after the obligatory topical subject NP or Pronoun. The VC is marked by the Change-of-Situation Marker \textit{paa} \textsc{tm}3 or unmarked. The subject may be formed by a NP, a Heavy Pronoun (\ref{ex:sketch-74}), or a shortened proclitic form, e.g.  \textit{a=re} `then we' (1\textsc{pl}.\textsc{in}.\textsc{pron}=\textsc{consec}), \textit{o=re} `then he/she/it/they' (3.\textsc{pron}=\textsc{consec}), which exclusively occur in this clause type, see (\ref{ex:sketch-78}). 

\ea\label{ex:sketch-78}
    are paa varoava e, ore raka.\\
    \gll a=re	paa	va-roava	e	o=re	raka\\
    1\textsc{pl}.\textsc{in}.\textsc{pron}=\textsc{consec}	\textsc{tm}3	\textsc{caus}-be.in.the.sun	3\textsc{sg}.\textsc{pron}	3.\textsc{pron}=\textsc{consec}	dry\\
    \glt `(After having scraped the wild banana plant,) we put it in the sun so that it becomes dry.' (\ref{chap:11}:\ref{ex:text7-56})
\z

The construction \textit{ore} \textit{suguna} `until' (lit. `then it reaches') is a common way of expressing that some situation comes to an end, see (\ref{chap:5}:\ref{ex:text01-164}), (\ref{chap:34}:\ref{ex:text27-21}), (\ref{chap:37}:\ref{ex:text30-16}).\footnote{See: \url{https://dictionaria.clld.org/units/teop-suguna}.} 

\section{Alienable and Inalienable Possessive Constructions}\label{sec:posscon}

The Alienable Possessive Construction (Al.Poss.Constr) is formed by a NP that is headed by a Noun or a Verb and followed by a PrepP introduced by \textit{te} \textsc{prep}. This PrepP refers to the owner of the referent of the NP, e.g. \textit{a inu te=ara} `our house' (\ref{ex:sketch-22}), its consumer, e.g. \textit{a tabaan te=o burumakao} `the food of the cows' (\ref{ex:sketch-24}), and the actor, e.g. \textit{a kiu vahihiva te suburiori} `the difficult work of their grandparent' (\ref{ex:sketch-58}).

The Inalienable Possessive Construction (Inal.Poss.Constr.) is formed by a NP that with one exception  is marked by the Possessive Marker, see \tabref{tab:pronouns}, and is obligatorily followed by another NP, see (\ref{ex:sketch-79}),

\ea\label{ex:sketch-79}
    o paka no overe\\
    \gll o paka n=o overe\\
    \textsc{art}2.\textsc{pl} leaf 3\textsc{sg}.\textsc{poss}=\textsc{art}3.\textsc{sg} coconut.palm\\
    \glt `the leaves of the coconut palm' (\ref{chap:5}:\ref{ex:text01-15})
\z

\noindent or a Pronoun, see (\ref{ex:sketch-80}).

\ea\label{ex:sketch-80}
    A ba suburaara ...\\
    \gll a	ba	subu-ra-ara\\
    [\textsc{art}2.\textsc{sg}	\textsc{kin}.\textsc{plm}	grandparent-1\textsc{pl}.\textsc{in}.\textsc{poss}-1\textsc{pl}.\textsc{in}.\textsc{pron}]\textsubscript{NP}\\
    \glt `our grandparents ...' (\ref{chap:5}:\ref{ex:text01-75})
\z

The standard orthography of \textit{o paka no overe} `the leaves of the coconut palm' in (\ref{ex:sketch-79}) does not reflect its morphosyntactic structure, because the Possessive Marker \textit{n} 3\textsc{sg}.\textsc{poss} is suffixed to \textit{paka} `leaf' and signifies its relationality, whereas the Article \textit{o} introduces the complement NP \textit{o overe}. The Possessive Marker inflects for person and number and indexes the complement NP or Pronoun with the exception of the 1\textsuperscript{st} Person Singular Pronoun \textit{naa}, see \tabref{tab:pronouns} in \sectref{sec:fourcat}.

In (\ref{ex:sketch-80}) the segmentation of \textit{subu-ra-ara} `our grandparent' and its glossing render the morphological structure of the word, but in other examples and in the texts, the sequence of a Possessive Marker and a Basic Pronoun is treated as a single linguistic unit, e.g. \textit{a ba subu-raara} `our grandparents (\textsc{art}2.\textsc{sg} grandparent-1\textsc{pl}.\textsc{in}.\textsc{poss}.\textsc{pron}' (\ref{ex:sketch-1}), (\ref{ex:sketch-37}).

The Inal.Poss.Constr. denotes the following relationships:\footnote{Search (Words extra) in the \textsc{mted} for the Gloss ``POSS'' and find 137 entries.}

\begin{itemize}
    \item unalterable kinship, see (\ref{ex:sketch-80});
    \item a body part and the human or animal it belongs to (\ref{ex:sketch-33});
    \item a part of a whole (\ref{ex:sketch-25}); 
    \item the name of a person, animal, or thing (\ref{ex:sketch-51});
    \item the size of a thing (\ref{chap:5}:\ref{ex:text01-54}), (\ref{chap:16}:\ref{ex:text11-42}), (\ref{chap:19}:\ref{ex:text14-28});
    \item the quantity of things (\ref{chap:16}:\ref{ex:text11-29});
    \item a spatial relationship (\ref{ex:sketch-23}); (\ref{chap:5}:\ref{ex:text01-33}); (\ref{chap:5}:\ref{ex:text01-161}), (\ref{chap:32}:\ref{ex:text25-8});
    \item a temporal relationship (\ref{chap:22}:\ref{ex:text16-15}), (\ref{chap:22}:\ref{ex:text16-19}), 
    \item an action and its undergoer (\ref{ex:sketch-28});
    \item the relation between a property and the bearer of the property (\ref{ex:sketch-81});
    \item a comparative relationship.\footnote{See \textcite[279--280]{Mosel:2017as}, \url{https://dictionaria.clld.org/units/teop-babarii}, and \url{https://dictionaria.clld.org/units/teop-kirisi\_1}.}
\end{itemize}

The head of the following Inal.Poss.Constr. \textit{mataa} `good' is classified as an Adjective, because the selection of the Article \textit{a} (\textsc{art}2.\textsc{sg}) is determined by the preceding context, i.e. the surrendering which is \textit{a hata} `bad',  see (\ref{ex:sketch-81}).

\ea\label{ex:sketch-81}
    \ea\label{ex:sketch-81a}
        Beara hee komana ni eara, a hata,\\
        \gll be=ara hee komana ni e-ara a hata\\
        when=1\textsc{pl}.\textsc{in}.\textsc{pron} give self \textsc{appl} \textsc{h}-1\textsc{pl}.\textsc{in}.\textsc{pron}	\textsc{art}2.\textsc{sg} bad\\
        \glt `If we surrender (lit. give ourselves), it is bad,'
    \ex\label{ex:sketch-81b}
        evehee na tei me nana bona mataa nae mee.\\
        \gll evehee na tei me nana bona mataa nae mee\\
        but \textsc{tm}2 be.there with 3\textsc{sg}.\textsc{ipfv} \textsc{obj}.\textsc{art}2.\textsc{sg} good		3\textsc{sg}.\textsc{poss}.\textsc{pron} also\\
        \glt `but it also has its good side.' (\cite[105]{Mahaka:2014oc})
    \z
\z

Relational Nouns can be derelationalized by the suffix \textit{-na} `\textsc{derel}',\footnote{Search (Words extra) in the \textsc{mted} for the Gloss  ``DEREL'' and find 63 examples.} as, for example \textit{toon-} `top', see (\ref{ex:sketch-82}),

\ea\label{ex:sketch-82}
    bono kapa vaa toonuna,\\
    \gll bono	kapa	vaa	toonu-na\\
    \textsc{obj}.\textsc{art}3.\textsc{sg}	cover	\textsc{lnk}	top-\textsc{derel}\\
    \glt `(they start scraping off) the top cover (lit. the cover of the top),' (\ref{chap:8}:\ref{ex:text4-15})
\z

\noindent or the kinship term \textit{sina-} `mother', see (\ref{ex:sketch-83}).

\ea\label{ex:sketch-83}
    to sinivi kiri bona vuu tom sinana.\\
    \gll to	sinivi	ki=ri	bona	vuu	tom	sina=na\\
    \textsc{nspec}3.\textsc{sg}	canoe	\textsc{dat}=3\textsc{pl}.\textsc{objm}	\textsc{art}2.\textsc{sg}	cluster	\textsc{dyad}	mother=\textsc{derel}\\
    \glt `a canoe for the family.' (\ref{chap:7}:\ref{ex:text3-15})\footnote{Search (Words extra) in the \textsc{mted} for the Gloss  ``DYAD'' and find 11 examples.}
\z

Other nouns, e.g. \textit{paka} `leaf' are marked by the \textsc{art}2 Article, when they head a Inal.Poss.Constr. as in (\ref{ex:sketch-79}), but by the \textsc{art}3 Article (\ref{ex:sketch-69}), when they are used as a non-relational Noun.\footnote{Search the \textsc{mted} for the Part of Speech ``N2/N3'' and find 53 entries.; also see \textcite[59--64]{Mosel:2014kg}.} 

\section{Dependent Clauses}\label{sec:depclaus}
\subsection{Introduction}\label{sec:depclausintro}

There are two kinds of dependent clauses in Teop that are marked by a conjunction:

\begin{enumerate}
    \item Subordinate Clauses that hold the position of a clause constituent, 
    \item Dependent Non-Embedded Clauses. 
\end{enumerate}

Thirdly, there are unmarked dependent clauses in which a syntactically obligatory subject or object NP or Pronoun is missing, but understood from the context. For examples see \tabref{tab:zero} in \sectref{sec:fourcat}.

\subsection{Subordinate Clauses}\label{sec:subordclaus}

Subordinate clauses are found with 7 syntactic functions, see \tabref{tab:subclause}.


\begin{table}
    \caption{Subordinate clauses.} \label{tab:subclause}
\begin{tabularx}{\textwidth}{QXXXXl}
    \lsptoprule
    \textsc{functions}&\multicolumn{5}{c}{\textsc{subordinate clauses}}\\
    \cmidrule{2-6}
    &&\textit{be}\textsuperscript{1}-&\textit{be}\textsuperscript{2}-&\textit{tea-}&\textit{mene-}\\
    &\textsc{relative}&\textsc{adverbial}&\textsc{compl.}&\textsc{compl.}&\textsc{purposive}\\
    \midrule
        predicate&(\ref{chap:5}:\ref{ex:text01-56})&-&(\ref{chap:5}:\ref{ex:text01-176})&(\ref{chap:31}:\ref{ex:text24-8})&(\ref{chap:30}:\ref{ex:text23-6})\\
        subject&(\ref{chap:33}:\ref{ex:text26-11})&-&-&(\ref{chap:5}:\ref{ex:text01-146})&-\\
        object&-&-&(\ref{chap:5}:\ref{ex:text01-12})&(\ref{chap:5}:\ref{ex:text01-11})&-\\
        oblique object&\multirow{2}{*}{(\ref{chap:17}:\ref{ex:text12-40})}&\multirow{2}{*}{-}&\multirow{2}{*}{(\ref{chap:5}:\ref{ex:text01-164})}&\multirow{2}{*}{(\ref{chap:5}:\ref{ex:text01-177})}&\multirow{2}{*}{-}\\
        attribute&(\ref{chap:5}:\ref{ex:text01-57})&-&-&-&(\ref{chap:12}:\ref{ex:text8-7})\\
        compl. of \textsc{lnk}&\multirow{2}{*}{-}&\multirow{2}{*}{-}&\multirow{2}{*}{-}&\multirow{2}{*}{(\ref{chap:5}:\ref{ex:text01-180})}&-\\
        adjunct&(\ref{chap:17}:\ref{ex:text12-30})&(\ref{chap:5}:\ref{ex:text01-2})&-&-&(\ref{chap:29}:\ref{ex:text22-31})\\
    \lspbottomrule
\end{tabularx}
\end{table}


\noindent Comments:

\begin{itemize}
    \item The Relative Clause is introduced by the Relative Pronoun \textit{to} \textsc{rel}, unless the following VC starts with the TM-Marker \textit{tau} `about to' (\textsc{tm}5) or \textit{toro} `must' (\textsc{tm}7). The Relative Pronoun does not inflect and functions as a subject, object, oblique object or adjunct of the Relative clause. The Relative Clause functions as an attribute, and similar to NPs and PrepPs, as a predicate, subject, object, oblique object, and a locative or temporal Adjunct.
    \item Complement clauses introduced by \textit{tea} (\textsc{compl}1) are similar to English to-infinitive clauses. Either the subject or the object NP or Pronoun of the matrix clause functions as their subject. Those introduced by \textit{be}\textsuperscript{2} (\textsc{compl}2) have a subject NP that is different from the subject of the main clause. 
    \item The Purposive Clause is introduced by \textit{mene} `for' and structurally similar to the \textit{tea}-Complement Clause, but differs with respect to its syntactic functions and meaning, see \tabref{tab:subclause}.
    \item The Adverbial Clause is introduced by \textit{be}\textsuperscript{1} `when, if'. In the texts about the growth of plants and the production of artefacts temporal Adverbial Clauses typically precede a matrix clause that is frequently marked by the consecutive conjunction \textit{re} `then' (\textsc{consec}).
\end{itemize}


\subsection{Dependent Non-Embedded Clauses}\label{sec:depnonclause}

The three types of dependent non-embedded clauses are linked to a preceding clause by a conjunction:

\begin{enumerate}
    \item the Consecutive Clause (436 tokens) is marked by \textit{re} `then, so that' (\textsc{consec}) after a topical NP or Pronoun with, see \sectref{sec:consecpron} and \url{https://dictionaria.clld.org/units/teop-re}.
    \item the Adversative Clause (35 tokens) is marked by \textit{he} `but, while' (\textsc{advers}) is the same position as the conjunction re in the Consecutive Clause, see \url{https://dictionaria.clld.org/units/teop-he}. 
    \item the Coordinative Clause (19 tokens) is introduced by the conjunction \textit{me} `and' (and4), see \url{https://dictionaria.clld.org/units/teop-me\_1}.
\end{enumerate}

\section{Filtered Searches in the Multifunctional Teop-English Dictionary (\textsc{mted})}\label{sec:mted}

The \textsc{mted} has two interfaces for filtered searches. The first is called ``Words''\footnote{See: \url{https://dictionaria.clld.org/contributions/teop\#twords}.}, the second is called ``Words extra''\footnote{See: \url{https://dictionaria.clld.org/contributions/teop\#twords2}.}, see Figures \ref{fig:mtedwords} and \ref{fig:mtedwordsextra}. The column label ``MWE Structure'' stands for ``Multi-Word-Expression Structure''. Example searches are shown in Tables \ref{tab:mtedwordsearch} and \ref{tab:mtedwordextrasearch}.

\begin{figure}[H]
    \centering
    \includegraphics[width=1\linewidth]{figures/mted-1.png}
    \caption{Search interface ``Words''.}
    \label{fig:mtedwords}
\end{figure}


\begin{table}
    \caption{Examples for searches in two fields in ``Words''} \label{tab:mtedwordsearch}
\begin{tabularx}{\textwidth}{Xp{3cm}p{3cm}l}
    \lsptoprule
    \textsc{topic}&\multicolumn{2}{c}{\textsc{search query}}&\textsc{hits}\\
    &\textsc{field: ``keyword''}&\textsc{field: ``keyword''}&\\    \midrule
%     \endfirsthead
%     \multicolumn{4}{c}%
%     {\tablename\ \thetable{} -- continued from previous page}\\
%     \lsptoprule
%     \textsc{topic}&\multicolumn{2}{c}{\textsc{search query}}&\textsc{hits}\\
%     &\textsc{field: ``keyword''}&\textsc{field: ``keyword''}&\\	    \midrule
%     \endhead
%     \hline \multicolumn{4}{r}{{Continued on next page}}\\
%     \endfoot
%     \endlastfoot
        \multirow{2}{*}{Temporal adverbs}&Part of Speech: ``adv''&Semantic Domain: ``time''&\multirow{2}{*}{35}\\
        Intransitive VCs with a transitive verb&Part of Speech: ``vi.constr''&MWE Structure: ``vt''&\multirow{2}{*}{72}\\
        Transitive VCs with an incorporated preposition&Part of Speech: ``vt.constr''&MWE Structure: ``Prep''&\multirow{2}{*}{166}\\
    \lspbottomrule
\end{tabularx}
\end{table}


\begin{figure}[H]
    \centering
    \includegraphics[width=1\linewidth]{figures/mted-2.png}
    \caption{Search interface ``Words extra''.}
    \label{fig:mtedwordsextra}
\end{figure}


\begin{table}
    \caption{Examples for searches in two fields in ``Words extra''.} \label{tab:mtedwordextrasearch}
\begin{tabularx}{\textwidth}{Xp{3cm}p{3cm}l}
    \lsptoprule
    \textsc{topic}&\multicolumn{2}{c}{\textsc{search query}}&\textsc{hits}\\
    &\textsc{field: ``keyword''}&\textsc{field: ``keyword''}&\\
    \midrule
        \multirow{6}{*}{Possessive Constructions}&Part of Speech: ``n1.constr''&\multirow{2}{*}{Gloss: ``poss''}&\multirow{2}{*}{21}\\
        &Part of Speech: ``n2.constr''&\multirow{2}{*}{Gloss: ``poss''}&\multirow{2}{*}{86}\\
        &Part of Speech: ``n3.constr''&\multirow{2}{*}{Gloss: ``poss''}&\multirow{2}{*}{7}\\
        Intransitive Verbs with a causative prefix&\multirow{2}{*}{Part of Speech: ``vi''}&\multirow{2}{*}{Gloss: ``{\^{}}CAUS''}&\multirow{2}{*}{17}\\
        Transitive Verbs with a causative prefix&\multirow{2}{*}{Part of Speech: ``vt''}&\multirow{2}{*}{Gloss: ``{\^{}}CAUS''}&\multirow{2}{*}{192}\\
        NPs headed by a causative verb&Part of Speech: ``n2.constr''&\multirow{2}{*}{Gloss: ``{\^{}}CAUS''}&\multirow{2}{*}{2}\\
    \lspbottomrule
\end{tabularx}
\end{table}


\end{document}
